% =====================================================================
%  Rapport de projet — EquiPay
%  Compilation : latexmk -pdf main.tex
%  (ou pdflatex → bibtex → pdflatex → pdflatex)
% =====================================================================
\documentclass[12pt,letterpaper]{article}

% ---------- Langue et polices ----------
\usepackage[T1]{fontenc}
\usepackage[utf8]{inputenc}
\usepackage[french]{babel}
\usepackage{lmodern}
\usepackage{textcomp}
\usepackage{microtype}

% ---------- Mise en page ----------
\usepackage[left=1in,right=1in,top=1in,bottom=1in,footskip=0.5in]{geometry}
\usepackage{setspace}
\setstretch{1.35}
\setlength{\parindent}{1.5em}
\setlength{\parskip}{0.35em}

% ---------- Couleurs (reprises du rapport de référence et de l'app) ----------
\usepackage[table]{xcolor}
\definecolor{bleu}{HTML}{1A365D}      % titres
\definecolor{vert}{HTML}{3A9D5D}      % bande EMSI, filets de chapitre
\definecolor{bordeaux}{HTML}{990033}  % carré de la page de garde
\definecolor{lagon}{HTML}{0D9488}     % couleur principale d'EquiPay
\definecolor{lagonclair}{HTML}{CCFBF1}
\definecolor{orange}{HTML}{F97316}    % accent d'EquiPay
\definecolor{orangeclair}{HTML}{FFEDD5}
\definecolor{gristexte}{HTML}{64748B}
\definecolor{grisfond}{HTML}{F1F5F9}
\definecolor{entete}{HTML}{E8EEF6}    % en-tête des tableaux
\definecolor{recevoir}{HTML}{16A34A}
\definecolor{rembourser}{HTML}{DC2626}
\definecolor{violet}{HTML}{8B5CF6}
\definecolor{bleuclair}{HTML}{3B82F6}
\definecolor{ambre}{HTML}{F59E0B}
\definecolor{codefond}{HTML}{1E1E2E}
\definecolor{codetexte}{HTML}{E2E8F0}
\definecolor{codemot}{HTML}{C792EA}
\definecolor{codechaine}{HTML}{C3E88D}
\definecolor{codecomm}{HTML}{7F8C98}

% ---------- Figures, tableaux, listes ----------
\usepackage{graphicx}
\graphicspath{{images/}}
\usepackage{float}
\usepackage{caption}
\usepackage{subcaption}
\usepackage{array}
\usepackage{tabularx}
\usepackage{multirow}
\usepackage{enumitem}
\setlist[itemize]{label=---, leftmargin=2.2em, itemsep=0.2em, topsep=0.4em}
\setlist[enumerate]{leftmargin=2.2em, itemsep=0.2em, topsep=0.4em}
\renewcommand{\arraystretch}{1.45}
\newcolumntype{L}[1]{>{\raggedright\arraybackslash}p{#1}}
\newcolumntype{Y}{>{\raggedright\arraybackslash}X}
\newcolumntype{C}[1]{>{\centering\arraybackslash}p{#1}}

% Légendes : « Figure 3 : ... », étiquette en petites capitales bleues
\captionsetup{font=small, labelfont={sc,color=bleu}, labelsep=colon, justification=centering, skip=6pt}
\captionsetup[sub]{font=footnotesize, labelfont={}, labelformat=parens, skip=3pt}
\addto\captionsfrench{\renewcommand{\tablename}{Tableau}}

% ---------- Dessins ----------
\usepackage{tikz}
\usetikzlibrary{arrows.meta,positioning,shapes.geometric,shapes.multipart,calc,fit,backgrounds,shadows}
\usepackage{pgfplots}
\pgfplotsset{compat=1.18}
\usepackage{pgf-umlsd}

% ---------- Code ----------
\usepackage{listings}
\lstset{
  backgroundcolor=\color{codefond},
  basicstyle=\ttfamily\scriptsize\color{codetexte},
  keywordstyle=\color{codemot}\bfseries,
  stringstyle=\color{codechaine},
  commentstyle=\color{codecomm}\itshape,
  numbers=none,
  breaklines=true,
  breakatwhitespace=true,
  columns=fullflexible,
  keepspaces=true,
  showstringspaces=false,
  frame=single, framerule=0pt, framesep=8pt,
  xleftmargin=8pt, xrightmargin=8pt,
  aboveskip=4pt, belowskip=4pt,
  upquote=true,
  literate=
    {é}{{\'e}}1 {è}{{\`e}}1 {ê}{{\^e}}1 {ë}{{\"e}}1 {à}{{\`a}}1 {â}{{\^a}}1
    {ç}{{\c{c}}}1 {ô}{{\^o}}1 {î}{{\^i}}1 {ï}{{\"i}}1 {ù}{{\`u}}1 {û}{{\^u}}1
    {É}{{\'E}}1 {À}{{\`A}}1 {«}{{\guillemotleft}}1 {»}{{\guillemotright}}1
    {’}{{'}}1 {→}{{$\rightarrow$}}2 {…}{{\ldots}}3 {−}{{-}}1
}
\lstdefinelanguage{Dart}{
  morekeywords={class,final,const,var,for,if,while,return,List,double,int,void,new,in,true,false,null,this,import,as},
  sensitive=true, morecomment=[l]{//}, morestring=[b]', morestring=[b]"
}
\lstdefinelanguage{SQL2}{
  morekeywords={create,policy,on,for,select,insert,update,delete,to,using,with,check,and,or,exists,from,where,in,returns,language,stable,as,function,replace,alter,table,enable,row,level,security,authenticated,set},
  sensitive=false, morecomment=[l]{--}, morestring=[b]'
}

% ---------- Titres ----------
\usepackage{titlesec}
% Chapitres (niveau \section) : « 1. Titre » en bleu, filet vert dessous
\titleformat{\section}{\Large\color{bleu}}{\thesection.}{0.5em}{}[{\vspace{2pt}\color{vert}\titlerule[0.9pt]}]
\titlespacing*{\section}{0pt}{0pt}{1.2em}
\titleformat{\subsection}{\large\bfseries\color{bleu}}{\thesubsection}{0.8em}{}
\titlespacing*{\subsection}{0pt}{1.4em}{0.6em}
\titleformat{\subsubsection}{\normalsize\bfseries\itshape\color{bleu}}{\thesubsubsection}{0.8em}{}
\titlespacing*{\subsubsection}{0pt}{1.1em}{0.4em}
% Chaque chapitre commence sur une nouvelle page
\newcommand{\chapitre}[1]{\clearpage\section{#1}}

% Titres des pages liminaires (REMERCIEMENTS, RÉSUMÉ...) : centrés, filet dessous, entrée dans la table
\newcommand{\pagetitre}[1]{%
  \clearpage
  \phantomsection
  \addcontentsline{toc}{liminaire}{#1}%
  \begin{center}
    {\Large\bfseries #1}\par\vspace{1.4em}
    \rule{\textwidth}{0.6pt}
  \end{center}
  \vspace{1.2em}%
}

% ---------- Table des matières et listes ----------
\usepackage{tocloft}
\setcounter{tocdepth}{2}
\addto\captionsfrench{%
  \renewcommand{\contentsname}{TABLE DES MATIÈRES}%
  \renewcommand{\listfigurename}{LISTE DES FIGURES}%
  \renewcommand{\listtablename}{LISTE DES TABLEAUX}}
% Même présentation que les autres pages liminaires : titre centré, filet dessous
\newcommand{\filetliste}{\hfill\null\par\nobreak\vspace{1.1em}\noindent\rule{\textwidth}{0.6pt}\par\vspace{0.4em}}
\renewcommand{\cfttoctitlefont}{\hfill\Large\bfseries}
\renewcommand{\cftaftertoctitle}{\filetliste}
\renewcommand{\cftloftitlefont}{\hfill\Large\bfseries}
\renewcommand{\cftafterloftitle}{\filetliste}
\renewcommand{\cftlottitlefont}{\hfill\Large\bfseries}
\renewcommand{\cftafterlottitle}{\filetliste}
\setlength{\cftbeforetoctitleskip}{0pt}
\setlength{\cftbeforeloftitleskip}{0pt}
\setlength{\cftbeforelottitleskip}{0pt}
\renewcommand{\cftsecfont}{\bfseries}
\renewcommand{\cftsecpagefont}{\bfseries}
\renewcommand{\cftsecaftersnum}{.}
\setlength{\cftbeforesecskip}{0.9em}
\setlength{\cftbeforesubsecskip}{0.15em}
\setlength{\cftsubsecindent}{0pt}
\setlength{\cftsubsecnumwidth}{2.6em}
\setlength{\cftfignumwidth}{2.4em}
\setlength{\cfttabnumwidth}{2.4em}
\renewcommand{\cftfigpresnum}{}
\renewcommand{\cfttabpresnum}{}
% Niveau « liminaire » : pages hors chapitres, en gras et sans numéro
\makeatletter
\newcommand*{\l@liminaire}[2]{\par\vspace{0.75em}\noindent{\bfseries #1}\hfill{\bfseries #2}\par}
\makeatother
\setlength{\cftaftertoctitleskip}{0pt}
\setlength{\cftafterloftitleskip}{0pt}
\setlength{\cftafterlottitleskip}{0pt}

% ---------- En-têtes et pieds de page ----------
\usepackage{fancyhdr}
\pagestyle{fancy}
\fancyhf{}
\fancyfoot[C]{\small\thepage}
\renewcommand{\headrulewidth}{0pt}
\fancypagestyle{plain}{\fancyhf{}\fancyfoot[C]{\small\thepage}\renewcommand{\headrulewidth}{0pt}}

% ---------- Encadrés ----------
\usepackage[most]{tcolorbox}
\newtcolorbox{encadre}[1][]{enhanced, colback=lagonclair!45, colframe=lagon, boxrule=0pt,
  leftrule=3pt, arc=2pt, left=10pt, right=10pt, top=6pt, bottom=6pt, fontupper=\small, #1}

% ---------- Divers ----------
\usepackage{csquotes}
\usepackage{url}
% Bibliographie (BibTeX, style unsrt-fr) : numérotée dans l'ordre des citations,
% seules les sources citées apparaissent
\usepackage[numbers,sort&compress]{natbib}
\setlength{\bibsep}{0.55em}
\renewcommand{\bibfont}{\small}
\renewcommand{\bibsection}{\pagetitre{RÉFÉRENCES}}
\usepackage[hidelinks]{hyperref}
\hypersetup{pdftitle={EquiPay - Rapport de projet}, pdfauthor={Alae Lahbichi, Haytham Benladoul}}
\usepackage{xspace}
\newcommand{\equipay}{EquiPay\xspace}
\newcommand{\code}[1]{\texttt{\small #1}}
\newcommand{\ecran}[2]{\includegraphics[width=#1\textwidth]{#2}}
% Capture de téléphone : bordure fine et coins nets
\newcommand{\capture}[2][0.3]{\fcolorbox{gristexte!40}{white}{\includegraphics[width=#1\textwidth]{#2}}}

\begin{document}

% =====================================================================
%  Page de garde (sans numéro, non comptée)
% =====================================================================
% ---------------------------------------------------------------------
%  Page de garde : aucun numéro affiché, page non comptée
% ---------------------------------------------------------------------
\begin{titlepage}
\thispagestyle{empty}
\setstretch{1}
\begin{tikzpicture}[remember picture, overlay]
  % Bande verte à gauche et carré bordeaux en bas
  \fill[vert] ([xshift=0.58in, yshift=-0.58in]current page.north west) rectangle ([xshift=0.9in, yshift=0.82in]current page.south west);
  \fill[bordeaux] ([xshift=0.58in, yshift=0.62in]current page.south west) rectangle ([xshift=0.9in, yshift=0.3in]current page.south west);
  % Logos : école à gauche, application à droite
  \node[anchor=north west, inner sep=0] at ([xshift=1.17in, yshift=-1.12in]current page.north west)
    {\includegraphics[width=2.75in]{logo_emsi.png}};
  \node[anchor=north east, inner sep=0] at ([xshift=-0.85in, yshift=-0.9in]current page.north east)
    {\includegraphics[width=1.05in]{logo_equipay.png}};
\end{tikzpicture}

\begin{center}
  \vspace*{1.55in}
  {\Large\color{bleu} Rapport de Projet}\\[1.1em]
  {\large\bfseries\color{bleu} 5\textsuperscript{ème} année}\\[0.15em]
  {\large\bfseries\color{bleu} IASDR -- Intelligence Artificielle}\\[2.2em]
  {\normalsize\bfseries Sous le thème}\\[1.4em]
  \rule{0.62\textwidth}{0.6pt}\\[1.4em]
  {\fontsize{21}{27}\selectfont EquiPay : application intelligente avec\\ un assistant conversationnel et une\\ automatisation basées sur l'IA\par}
  \vspace{1.3em}
  \rule{0.62\textwidth}{0.6pt}
\end{center}

\vspace{3em}
\noindent
\begin{minipage}[t]{0.6\textwidth}
  {\bfseries\color{bleu} Réalisé par :}\\
  Alae Lahbichi\\
  Haytham Benladoul\\[0.9em]
  {\bfseries\color{bleu} Encadré par :}\\
  Hasnââ CHAABI
\end{minipage}

\vfill
\begin{center}
  {\bfseries\color{bleu} Année Universitaire : 2026-2027}
\end{center}
\vspace*{0.1in}
\end{titlepage}


% =====================================================================
%  Pages liminaires : numérotation en chiffres romains
% =====================================================================
\pagenumbering{roman}
\setcounter{page}{1}
% ---------------------------------------------------------------------
%  Pages liminaires (numérotées en chiffres romains)
% ---------------------------------------------------------------------

\pagetitre{REMERCIEMENTS}

Nous tenons tout d'abord à remercier notre encadrante, Madame Hasnââ CHAABI, pour
la qualité de son accompagnement tout au long de ce projet. Ses conseils, sa
disponibilité et ses remarques nous ont aidés à garder une démarche rigoureuse, en
particulier sur l'évaluation des modèles et sur la façon de présenter nos résultats.

Nous remercions également l'École Marocaine des Sciences de l'Ingénieur (EMSI) pour la
formation reçue et pour les moyens mis à notre disposition durant notre cursus
d'ingénieur, en 5\textsuperscript{ème} année de la filière IASDR, option Intelligence
Artificielle. Nos remerciements vont aussi à l'ensemble des enseignants qui ont
contribué à notre formation, notamment en traitement automatique du langage, en
apprentissage profond, en génération de texte et en développement logiciel : ce projet
s'appuie directement sur ce qu'ils nous ont transmis.

Nous remercions enfin nos familles et nos proches pour leur soutien et leurs
encouragements, ainsi que les amis qui ont accepté de tester l'application et de nous
faire part de leurs retours. Ce projet a été pour nous l'occasion de construire, de bout
en bout, une application réelle qui associe développement mobile et intelligence
artificielle, et de mesurer ce qu'il faut pour faire passer un modèle de l'entraînement
à un usage quotidien.

% ---------------------------------------------------------------------
\pagetitre{RÉSUMÉ}

Ce projet porte sur la conception et le développement d'\equipay, une application
mobile de partage des dépenses de groupe (voyages, sorties, colocation, crédits) qui
intègre plusieurs modèles d'intelligence artificielle. L'objectif est double : permettre
à un groupe de savoir à tout moment qui doit combien à qui, et supprimer la saisie
fastidieuse qui décourage les utilisateurs des applications existantes.

L'application, développée avec Flutter, s'appuie sur Supabase (PostgreSQL, authentification,
stockage et temps réel) pour les données partagées, et sur un backend FastAPI exécuté
localement sur un Mac Apple Silicon pour les fonctions d'IA. Quatre usages de l'IA
générative y sont intégrés : un assistant conversationnel qui répond aux questions sur
l'application, un scanner de tickets de caisse fondé sur un modèle vision-langage, la
saisie d'une dépense à partir d'une simple phrase, et l'organisation automatique du
programme d'un voyage. Les trois premiers modèles sont des versions affinées par LoRA de
Qwen3-4B et Qwen3-VL-4B, entraînées sur GPU puis converties en 4 bits avec MLX pour
l'inférence ; le quatrième combine un calcul d'itinéraire (OSRM), un algorithme d'ordre de
visite et un modèle de langage chargé de rédiger le programme. Chaque entraînement est
suivi dans MLflow, et seule la meilleure version, désignée « champion », est servie par
l'application. Sur leurs jeux de test respectifs, l'assistant passe d'un score de 0,276 à
0,770, le scanner de 0,792 à 0,954 et la saisie par phrase de 0,656 à 0,976.

\noindent\textbf{Mots-clés :} Intelligence artificielle, Génération de texte, LLM,
Modèle vision-langage, LoRA, Quantification, MLX, MLOps, MLflow, Flutter, Supabase,
FastAPI, Partage de dépenses.

% ---------------------------------------------------------------------
\pagetitre{ABSTRACT}

This project covers the design and development of \equipay, a mobile application for
sharing group expenses (trips, outings, shared housing, loans) that embeds several
artificial intelligence models. The goal is twofold: let a group know at any time who
owes how much to whom, and remove the tedious data entry that discourages users of
existing apps.

The application is built with Flutter and relies on Supabase (PostgreSQL, authentication,
storage and real-time updates) for shared data, and on a FastAPI backend running locally
on an Apple Silicon Mac for AI features. Four generative AI use cases are integrated: a
conversational assistant that answers questions about the app, a receipt scanner based on
a vision-language model, expense entry from a single free-text sentence, and automatic
planning of a trip itinerary. The first three models are LoRA fine-tuned versions of
Qwen3-4B and Qwen3-VL-4B, trained on a GPU and then quantized to 4 bits with MLX for
inference; the fourth combines route computation (OSRM), a visit-ordering algorithm and a
language model that writes the final plan. Every training run is tracked in MLflow, and
only the best version, tagged as ``champion'', is served to the app. On their respective
test sets, the assistant improves from a score of 0.276 to 0.770, the receipt scanner from
0.792 to 0.954, and sentence-based entry from 0.656 to 0.976.

\noindent\textbf{Keywords:} Artificial Intelligence, Text Generation, LLM,
Vision-Language Model, LoRA, Quantization, MLX, MLOps, MLflow, Flutter, Supabase,
FastAPI, Expense Sharing.

% ---------------------------------------------------------------------
%  Table des matières, liste des figures, liste des tableaux
% ---------------------------------------------------------------------
\clearpage
\phantomsection
\addcontentsline{toc}{liminaire}{TABLE DES MATIÈRES}
{\setstretch{1.15}\tableofcontents}

\clearpage
\phantomsection
\addcontentsline{toc}{liminaire}{LISTE DES FIGURES}
{\setstretch{1.15}\listoffigures}

\clearpage
\phantomsection
\addcontentsline{toc}{liminaire}{LISTE DES TABLEAUX}
{\setstretch{1.15}\listoftables}

% ---------------------------------------------------------------------
\pagetitre{LISTE DES ABRÉVIATIONS}

\begin{table}[H]
\centering
\small
\begin{tabularx}{0.92\textwidth}{|L{2.6cm}|Y|}
\hline
\rowcolor{entete}\textbf{Abréviation} & \textbf{Signification} \\ \hline
IA      & Intelligence Artificielle \\ \hline
LLM     & \textit{Large Language Model}, grand modèle de langage \\ \hline
VLM     & \textit{Vision-Language Model}, modèle vision-langage \\ \hline
LoRA    & \textit{Low-Rank Adaptation}, adaptation par matrices de faible rang \\ \hline
PEFT    & \textit{Parameter-Efficient Fine-Tuning}, affinage économe en paramètres \\ \hline
MLX     & Bibliothèque d'apprentissage automatique d'Apple pour Apple Silicon \\ \hline
MLOps   & \textit{Machine Learning Operations}, industrialisation des modèles \\ \hline
GPU     & \textit{Graphics Processing Unit}, processeur graphique \\ \hline
API     & \textit{Application Programming Interface}, interface de programmation \\ \hline
REST    & \textit{Representational State Transfer} \\ \hline
JSON    & \textit{JavaScript Object Notation} \\ \hline
JWT     & \textit{JSON Web Token}, jeton d'authentification signé \\ \hline
RLS     & \textit{Row Level Security}, sécurité au niveau des lignes \\ \hline
SQL     & \textit{Structured Query Language} \\ \hline
OCR     & \textit{Optical Character Recognition}, reconnaissance de caractères \\ \hline
OSRM    & \textit{Open Source Routing Machine}, moteur de calcul d'itinéraires \\ \hline
UML     & \textit{Unified Modeling Language} \\ \hline
UI / UX & Interface utilisateur / expérience utilisateur \\ \hline
MAD, DH & Dirham marocain \\ \hline
\end{tabularx}
\caption{Liste des abréviations utilisées dans le rapport}
\end{table}


% =====================================================================
%  Corps du rapport : numérotation en chiffres arabes à partir de 1
% =====================================================================
\clearpage
\pagenumbering{arabic}
\setcounter{page}{1}
% ---------------------------------------------------------------------
%  Introduction générale (première page en chiffres arabes)
% ---------------------------------------------------------------------
\pagetitre{INTRODUCTION GÉNÉRALE}

Partir en voyage entre amis, organiser une sortie ou partager un appartement amène
toujours la même question : qui a payé quoi, et qui doit combien à qui ? Les dépenses sont
avancées par des personnes différentes, parfois dans plusieurs devises, et rarement
réparties à parts égales. Après quelques jours, les comptes deviennent flous, et le calcul
final finit souvent sur un coin de feuille ou dans un tableur partagé. Des applications
comme Tricount~\cite{tricount} ou Splitwise~\cite{splitwise} ont simplifié ce calcul, mais elles reposent toujours sur une
saisie manuelle de chaque dépense, et n'aident en rien à préparer le voyage lui-même.

Dans le même temps, les progrès des grands modèles de langage ouvrent de nouvelles
possibilités. Un ticket de caisse, une phrase comme « J'ai payé 45 € de taxi pour Sara et
moi » ou une liste de lieux à visiter contiennent déjà toute l'information utile : un
modèle bien entraîné peut l'extraire et la transformer en données structurées, tout en
restant assez compact pour être affiné sur un GPU grand public et exécuté sur un portable.

C'est dans ce contexte que s'inscrit ce projet, réalisé dans le cadre du module de
Génération Textuelle de la 5\textsuperscript{ème} année à l'EMSI. Il vise à concevoir et à
développer \equipay, une application mobile de partage des dépenses de groupe qui utilise
l'intelligence artificielle là où elle apporte un gain concret : un assistant
conversationnel pour accompagner l'utilisateur, un scanner de tickets, la saisie d'une
dépense en une phrase et l'organisation automatique du programme d'un voyage. L'application
repose sur Flutter pour le mobile, sur Supabase pour les comptes et les données partagées,
et sur un backend FastAPI qui sert les modèles en local. Ces modèles sont des versions
affinées de Qwen3, suivies dans MLflow depuis la préparation des données jusqu'au choix de
la version mise en production.

Ce rapport est organisé en cinq chapitres. Le premier présente le cadre du projet et la
démarche suivie, le deuxième l'étude préalable et la spécification des besoins, et le
troisième la conception de l'application. Le quatrième détaille la réalisation, en
particulier l'entraînement et l'évaluation des modèles, et le dernier présente
l'application à travers ses principaux écrans.

% =====================================================================
\chapitre{Contexte général du projet}
% =====================================================================

Ce premier chapitre situe le projet. Il présente le cadre dans lequel il a été mené et
l'organisation de notre binôme, puis le besoin auquel répond \equipay et la problématique
qui en découle. Il fixe enfin les objectifs, le périmètre retenu et la démarche de travail
adoptée.

\subsection{Cadre du projet}

Le tableau~\ref{tab:cadre} résume le cadre dans lequel ce projet a été mené.

\begin{table}[H]
\centering
\small
\begin{tabularx}{\textwidth}{|L{3.2cm}|Y|}
\hline
\rowcolor{entete}\textbf{Critère} & \textbf{Description} \\ \hline
Établissement & École Marocaine des Sciences de l'Ingénieur (EMSI), 5\textsuperscript{ème} année du
cycle ingénieur, filière IASDR, option Intelligence Artificielle. \\ \hline
Module & Génération Textuelle : grands modèles de langage, affinage, évaluation de textes générés
et mise en production des modèles. \\ \hline
Équipe & Binôme : Alae Lahbichi et Haytham Benladoul. \\ \hline
Encadrement & Mme Hasnââ CHAABI. \\ \hline
Matériel & Un MacBook Pro M3 Pro (18 Go de mémoire unifiée) pour le développement, l'inférence et
le suivi des expériences ; un PC équipé d'une carte NVIDIA RTX 5080 pour l'entraînement des
modèles ; un iPhone pour les tests en conditions réelles. \\ \hline
Contrainte principale & Aucun service payant : services gratuits ou exécutés en local, modèles
ouverts. \\ \hline
\end{tabularx}
\caption{Cadre du projet}
\label{tab:cadre}
\end{table}

\subsection{Organisation du binôme}

Nous avons mené le projet de bout en bout, de l'idée de départ à l'application installée sur
téléphone. Le travail a été découpé par module (application mobile, base de données,
backend, puis chacun des modèles) afin que chacun puisse avancer sur sa partie sans bloquer
l'autre. Les choix d'architecture, les jeux de données et les résultats d'évaluation ont en
revanche toujours été discutés à deux, et chaque modification importante a été relue par
l'autre membre du binôme avant d'être intégrée. Les points réguliers avec notre encadrante
ont servi à valider les grandes orientations, en particulier la méthode d'évaluation des
modèles.

\subsection{Le partage des dépenses : un besoin du quotidien}

Les dépenses partagées font partie de la vie de groupe : un week-end, un voyage d'études,
une colocation ou un simple dîner. Le calcul n'est pas difficile en soi, mais il devient
vite pénible lorsque les dépenses s'accumulent, que les parts ne sont pas égales ou que
certains ont payé dans une autre devise. Pour un voyage, la question ne s'arrête d'ailleurs
pas à l'argent : il faut aussi décider de ce que l'on visite, dans quel ordre et à quel
moment, ce qui se fait le plus souvent dans une conversation de groupe où les propositions
se perdent.

\subsection{Problématique}

Deux difficultés se dégagent. La première tient à la saisie : chaque dépense doit être
entrée à la main (montant, payeur, participants, répartition), si bien que beaucoup
d'utilisateurs abandonnent en cours de route et reviennent au calcul approximatif. La
seconde tient à l'organisation : rien ne relie les dépenses du voyage à son programme, et
aucun outil grand public n'aide un groupe à choisir ensemble les lieux à visiter puis à les
enchaîner sans perdre de temps.

Les modèles de langage récents permettent de répondre aux deux. Un modèle vision-langage
peut lire un ticket de caisse, un modèle de langage affiné peut comprendre une phrase libre
et en tirer une dépense structurée, et un modèle généraliste peut rédiger un programme
clair à partir d'un itinéraire calculé. La question est alors de savoir comment intégrer
ces modèles dans une application fiable, gratuite et utilisable au quotidien, tout en
gardant la maîtrise de leur qualité : c'est la problématique à laquelle répond ce projet.

\subsection{Objectifs visés et délimitation du périmètre}

L'objectif est de concevoir une application mobile de partage des dépenses qui réduit la
saisie au minimum grâce à l'IA, aide le groupe à organiser ses sorties, et s'appuie sur une
démarche MLOps complète pour l'entraînement et le suivi des modèles.

Le périmètre retenu est présenté dans la figure~\ref{fig:perimetre}. Il comprend
l'application mobile (iOS en priorité, Android prêt), la base de données et
l'authentification, le backend d'IA, l'entraînement des modèles et leur suivi dans MLflow.
Le paiement intégré, la publication sur les magasins d'applications, l'hébergement du
backend dans le cloud et les notifications sont laissés comme évolutions futures.

\begin{figure}[H]
\centering
\begin{tikzpicture}[
  titre/.style={rounded corners=4pt, minimum width=6.6cm, minimum height=0.9cm, font=\bfseries\color{white}},
  item/.style={rounded corners=3pt, minimum width=6.6cm, minimum height=0.72cm, font=\footnotesize, draw=#1!60, fill=#1!8},
  ]
  \node[titre, fill=lagon] (in) at (0,0) {Inclus dans le projet};
  \node[titre, fill=gristexte] (out) at (7.3,0) {Hors périmètre (évolutions)};
  \foreach \t [count=\i] in {{Application mobile Flutter (iOS, Android prêt)},
                             {Comptes, activités, dépenses et soldes},
                             {Base Supabase : RLS et temps réel},
                             {Assistant conversationnel},
                             {Scanner de tickets et saisie par phrase},
                             {Programme de voyage et trajet optimisé},
                             {Entraînement et suivi dans MLflow}}
    \node[item=lagon] at (0,-0.35-0.85*\i) {\t};
  \foreach \t [count=\i] in {{Paiement intégré entre membres},
                             {Publication App Store et Play Store},
                             {Backend d'IA hébergé dans le cloud},
                             {Notifications push},
                             {Connexion avec un compte Apple}}
    \node[item=gristexte] at (7.3,-0.35-0.85*\i) {\t};
\end{tikzpicture}
\caption{Périmètre fonctionnel retenu : fonctionnalités incluses et exclues}
\label{fig:perimetre}
\end{figure}

\subsection{Démarche de travail et organisation du projet}

Nous avons adopté une démarche itérative et incrémentale, inspirée des méthodes agiles :
chaque cycle court se termine par une fonctionnalité testable sur téléphone. Comme le
montre la figure~\ref{fig:phases}, le projet s'est déroulé en cinq phases. Le cadrage a
fixé les besoins, le périmètre et les choix techniques. La conception a porté sur le schéma
de la base, l'architecture et la charte graphique. Le développement a produit l'application
et le backend. L'entraînement des modèles s'est déroulé en parallèle : préparation des
données, affinage sur GPU, évaluation et conversion. Les tests et la finalisation ont
regroupé la validation sur iPhone, l'optimisation et la rédaction de ce rapport.

\begin{figure}[H]
\centering
\begin{tikzpicture}[
  phase/.style={circle, minimum size=1.95cm, draw=white, line width=3pt, fill=#1, align=center,
                font=\footnotesize\bfseries, text=white, drop shadow={opacity=0.18}},
  bulle/.style={rectangle, rounded corners=6pt, draw=black!15, fill=white, align=center,
                font=\scriptsize, text width=2.75cm, minimum height=1.1cm},
  ]
  \fill[grisfond, rounded corners=6pt] (-1.6,-3.6) rectangle (14.1,3.6);
  \foreach \x/\y/\c/\n/\t/\d [count=\i] in {
      0/0.75/lagon/01/Cadrage/{Besoins, objectifs,\\ choix techniques},
      3.1/-0.75/violet/02/Conception/{Base de données,\\ architecture,\\ charte graphique},
      6.25/0.75/orange/03/Développement/{Application mobile,\\ backend, intégration\\ des modèles},
      9.4/-0.75/bleuclair/04/{Entraînement\\ des modèles}/{Données, affinage\\ LoRA, évaluation,\\ conversion MLX},
      12.5/0.75/recevoir/05/{Tests et\\ finalisation}/{Validation sur iPhone,\\ optimisation,\\ rapport}}
  {
    \node[phase=\c] (p\i) at (\x,\y) {\n\\[1pt]\t};
    \ifodd\i
      \node[bulle] (b\i) at (\x,-2.55) {\d};
      \draw[\c, thick] (p\i.south) -- (b\i.north);
      \fill[\c] (p\i.south) circle (2pt);
    \else
      \node[bulle] (b\i) at (\x,2.55) {\d};
      \draw[\c, thick] (p\i.north) -- (b\i.south);
      \fill[\c] (p\i.north) circle (2pt);
    \fi
  }
\end{tikzpicture}
\caption{Les cinq phases de la démarche de travail adoptée}
\label{fig:phases}
\end{figure}

\subsection{Conclusion du chapitre}

Ce chapitre a présenté le cadre du projet, l'organisation de notre binôme et le besoin
auquel répond \equipay. La problématique posée, les objectifs, le périmètre retenu et la
démarche adoptée constituent le point de départ de l'étude préalable présentée dans le
chapitre suivant.

% =====================================================================
\chapitre{Étude préalable et état de l'art}
% =====================================================================

Avant de concevoir l'application, il fallait comprendre comment les groupes gèrent
aujourd'hui leurs dépenses et ce que proposent les outils existants. Ce chapitre présente
cette analyse, les concepts d'intelligence artificielle mobilisés, l'étude de faisabilité,
puis les besoins fonctionnels et non fonctionnels qui ont guidé nos choix.

\subsection{Analyse de l'existant}

Avant de concevoir l'application, nous avons observé comment les groupes autour de nous
gèrent aujourd'hui leurs dépenses communes. Dans la plupart des cas, la gestion reste
artisanale : chacun note ce qu'il a payé dans son téléphone ou dans la conversation du
groupe, puis une personne fait le calcul à la fin, souvent dans un tableur. Quand une
application est utilisée, elle ne l'est généralement que par un ou deux membres du groupe,
qui saisissent tout pour les autres. Cette façon de faire présente plusieurs limites :

\begin{itemize}
  \item les dépenses oubliées ou saisies deux fois faussent le résultat final ;
  \item la répartition inégale (une personne qui ne boit pas, une chambre plus chère) est
        difficile à suivre sans erreur ;
  \item personne ne sait, pendant le voyage, où il en est réellement ;
  \item les tickets de caisse sont rarement conservés, alors qu'ils contiennent le montant
        exact, la date et l'enseigne ;
  \item le programme du voyage vit à part, dans une conversation où les propositions se
        perdent et où personne ne sait ce que les autres préfèrent.
\end{itemize}

\subsection{Recueil des besoins}

Ces observations, complétées par les échanges avec notre encadrante et par les retours
d'amis qui ont testé les premières versions, nous ont permis d'identifier les besoins
suivants :

\begin{itemize}
  \item partager une activité avec un groupe et voir les soldes de chacun en temps réel ;
  \item ajouter une dépense en quelques secondes, sans remplir un long formulaire ;
  \item savoir précisément qui doit rembourser qui, avec le moins de virements possible ;
  \item préparer le programme d'un voyage à plusieurs, en votant pour les propositions ;
  \item visualiser ce programme sur une carte et l'organiser pour perdre le moins de temps
        possible en déplacement ;
  \item obtenir de l'aide sur l'application sans quitter l'écran en cours.
\end{itemize}

\subsection{Étude des solutions concurrentes}

Nous avons étudié les applications les plus utilisées pour le partage des dépenses et pour
la préparation des voyages (tableau~\ref{tab:concurrents}).

\begin{table}[H]
\centering
\small
\begin{tabularx}{\textwidth}{|L{3.4cm}|Y|}
\hline
\rowcolor{entete}\textbf{Application} & \textbf{Constat} \\ \hline
Tricount~\cite{tricount} & Partage des dépenses simple et gratuit, groupes sans création de compte
obligatoire. La saisie reste entièrement manuelle et l'application ne propose aucune aide
à l'organisation du voyage. \\ \hline
Splitwise~\cite{splitwise} & Référence du partage de dépenses, avec simplification des dettes. Les
fonctions avancées, dont la lecture des tickets, sont réservées à l'abonnement payant ; aucune
saisie en langage naturel. \\ \hline
Settle Up~\cite{settleup} & Partage des dépenses multi-devises, plutôt orienté vers les groupes
réguliers (colocation, couple). Même logique de formulaire à remplir pour chaque dépense. \\ \hline
Wanderlog~\cite{wanderlog} & Planification de voyage sur carte avec liste des lieux et itinéraire.
Le suivi des dépenses y est secondaire et le choix collectif des lieux n'est pas au cœur de
l'outil. \\ \hline
\end{tabularx}
\caption{Panorama des solutions concurrentes étudiées}
\label{tab:concurrents}
\end{table}

Aucune de ces solutions ne réunit, dans un même outil, le partage des dépenses, une saisie
assistée par l'IA et l'organisation collective du voyage. C'est ce positionnement que
vise \equipay : une seule application pour payer, équilibrer les comptes et décider
ensemble de ce que l'on fait.

\subsection{Concepts clés}

Le tableau~\ref{tab:concepts} définit les principaux concepts mobilisés dans ce projet et
précise leur rôle dans l'application.

\begin{table}[H]
\centering
\small
\begin{tabularx}{\textwidth}{|L{3.4cm}|Y|}
\hline
\rowcolor{entete}\textbf{Concept} & \textbf{Définition et rôle dans le projet} \\ \hline
Grand modèle de langage (LLM) & Réseau de neurones de type \textit{Transformer} entraîné sur de vastes
corpus pour prédire le mot suivant, capable de comprendre et de produire du texte. Nous utilisons
Qwen3-4B-Instruct-2507~\cite{qwen3,qwen3_4b} pour l'assistant, la saisie par phrase et la rédaction
du programme de voyage. \\ \hline
Modèle vision-langage (VLM) & Modèle qui reçoit une image et du texte et répond en texte. Qwen3-VL-4B~\cite{qwen3vl}
lit directement la photo d'un ticket, sans étape d'OCR séparée, et renvoie ses champs au format JSON. \\ \hline
Affinage LoRA & Plutôt que de réentraîner tous les poids, LoRA~\cite{lora} ajoute à certaines couches
deux petites matrices de faible rang, seules apprises. Moins de 1 \% des paramètres sont entraînés,
ce qui rend l'affinage d'un modèle de 4 milliards de paramètres possible sur un seul GPU. \\ \hline
Quantification & Représentation des poids sur moins de bits (ici 4 bits au lieu de 16). La mémoire
nécessaire est divisée par trois à quatre, pour une perte de qualité mesurée et limitée. Elle est
réalisée avec MLX~\cite{mlx}, la bibliothèque d'Apple pour les puces Apple Silicon. \\ \hline
MLOps & Ensemble des pratiques qui rendent un modèle traçable et reproductible : données versionnées,
expériences enregistrées, comparaison des versions et mise en service contrôlée. Nous utilisons
MLflow~\cite{mlflow,mlflow_docs} et son registre de modèles. \\ \hline
Sécurité au niveau des lignes (RLS) & Règles écrites dans PostgreSQL~\cite{postgres_rls} qui filtrent
chaque requête selon l'utilisateur connecté. Un membre ne peut lire que les activités dont il fait
partie, même si l'application était détournée. \\ \hline
\end{tabularx}
\caption{Concepts clés mobilisés dans le projet}
\label{tab:concepts}
\end{table}

\subsection{Étude de faisabilité}

Le tableau~\ref{tab:faisabilite} résume les quatre volets de l'étude de faisabilité menée
avant le développement.

\begin{table}[H]
\centering
\small
\begin{tabularx}{\textwidth}{|L{3.0cm}|Y|}
\hline
\rowcolor{entete}\textbf{Volet} & \textbf{Constat} \\ \hline
Technique & Flutter, Supabase et FastAPI sont des technologies matures et bien documentées. Flutter
permet de viser iOS et Android avec un seul code, et Supabase fournit d'emblée l'authentification,
le stockage des fichiers et le temps réel. \\ \hline
Données & Des jeux publics de tickets annotés existent (CORD, SROIE, WildReceipt), mais aucun ne
couvre les devises qui nous intéressent (dirham, euro). Pour l'assistant et la saisie par phrase,
aucun jeu n'existe : il faut les construire nous-mêmes. \\ \hline
Matériel & Un modèle de 4 milliards de paramètres s'affine en LoRA sur une carte de 16 Go (RTX 5080).
En 4 bits, les trois modèles servis tiennent ensemble dans la mémoire d'un Mac de 18 Go. \\ \hline
Coûts & Aucun coût d'API : les modèles sont ouverts et exécutés en local, Supabase est utilisé dans
son offre gratuite, et les cartes et itinéraires reposent sur des services gratuits (Apple Plans,
OSRM). \\ \hline
\end{tabularx}
\caption{Synthèse de l'étude de faisabilité}
\label{tab:faisabilite}
\end{table}

\subsection{Spécification des besoins fonctionnels}

Les besoins fonctionnels retenus sont regroupés par domaine :

\begin{itemize}
  \item \textbf{Comptes et profil} : inscription par e-mail ou avec Google, profil avec photo,
        thème clair ou sombre, suppression du compte.
  \item \textbf{Activités} : créer une activité (voyage, sortie, quotidien, crédit) avec sa
        devise, rejoindre une activité par code avec validation du créateur, gérer les rôles
        (membre ou spectateur), archiver ou supprimer.
  \item \textbf{Dépenses} : ajouter une dépense avec montant, devise, date, payeur et photo du
        reçu ; la répartir à parts égales, en pourcentages, en montants ou en parts.
  \item \textbf{Soldes} : calculer le solde de chacun en temps réel, proposer le minimum de
        remboursements, afficher graphiques et statistiques.
  \item \textbf{Programme} : pour les voyages et les sorties, lister les lieux à visiter avec
        adresse, date et heure, voter pour ou contre, afficher le programme en liste, en
        calendrier et sur une carte.
  \item \textbf{Intelligence artificielle} : scanner un ticket, décrire une dépense en une
        phrase, organiser le trajet du programme, poser une question à l'assistant.
\end{itemize}

\subsection{Spécification des besoins non fonctionnels}

\begin{figure}[H]
\centering
\begin{tikzpicture}[
  num/.style={circle, fill=white, draw=black!10, minimum size=0.9cm, font=\bfseries, drop shadow={opacity=0.2}},
  titre/.style={font=\large\bfseries, anchor=south},
  texte/.style={font=\scriptsize, align=center, text width=3.9cm, anchor=north},
  lien/.style={draw=#1, line width=0.8pt},
  ]
  % Demi-couronne en six secteurs
  \foreach \i/\c in {0/lagon, 1/violet, 2/bleuclair, 3/orange, 4/recevoir, 5/ambre} {
    \pgfmathsetmacro{\a}{180-\i*30-1.5}
    \pgfmathsetmacro{\b}{180-(\i+1)*30+1.5}
    \fill[\c!85] (\a:1.4) arc[start angle=\a, end angle=\b, radius=1.4] -- (\b:2.9)
                 arc[start angle=\b, end angle=\a, radius=2.9] -- cycle;
    \pgfmathsetmacro{\m}{180-\i*30-15}
    \pgfmathtruncatemacro{\n}{\i+1}
    \node[num] (n\n) at (\m:2.9) {0\n};
  }
  % Libellés : trois à gauche, trois à droite
  \foreach \k/\x/\y/\c/\t/\d in {
      1/-6.1/0.15/lagon/Performance/{Réponse de l'IA en quelques secondes : modèles MLX 4 bits chargés au premier usage.},
      2/-6.1/2.35/violet/Fiabilité/{Repli sur PyTorch, temps estimés si OSRM ne répond pas, bandeau hors ligne.},
      3/-6.1/4.55/bleuclair/Sécurité/{Règles RLS sur toutes les tables, jeton JWT vérifié, reçus privés.},
      4/6.1/4.55/orange/Sobriété/{Aucun service payant ; backend limité à environ 5,8 Go de mémoire.},
      5/6.1/2.35/recevoir/Ergonomie/{Material 3, mode sombre ; l'IA pré-remplit, l'utilisateur valide.},
      6/6.1/0.15/ambre/Maintenabilité/{Migrations versionnées, 89 tests automatiques, modèles tracés dans MLflow.}}
  {
    \node[titre, text=\c!80!black] (t\k) at (\x,\y) {\t};
    \node[texte] at (t\k.south) {\d};
    \ifnum\k<4
      \draw[lien=\c] (t\k.east) ++(0.1,-0.05) -- (n\k);
    \else
      \draw[lien=\c] (t\k.west) ++(-0.1,-0.05) -- (n\k);
    \fi
  }
\end{tikzpicture}
\caption{Synthèse des exigences non fonctionnelles retenues pour l'application}
\label{fig:non_fonctionnels}
\end{figure}

La figure~\ref{fig:non_fonctionnels} synthétise les six exigences non fonctionnelles
retenues : la performance (réponse rapide des modèles malgré une exécution locale), la
fiabilité (l'application continue de fonctionner quand un service externe ne répond pas),
la sécurité (chaque utilisateur ne voit que ses propres activités), la sobriété (aucun coût
et une empreinte mémoire maîtrisée), l'ergonomie (prise en main immédiate) et la
maintenabilité (code testé et modèles traçables).

\subsection{Synthèse et choix retenus}

L'étude de l'existant a montré qu'aucun outil ne combine partage des dépenses, saisie
assistée par l'IA et organisation collective du voyage. Sur cette base, nous avons retenu
les choix résumés dans la figure~\ref{fig:choix}.

\begin{figure}[H]
\centering
\begin{tikzpicture}[
  centre/.style={circle, fill=lagon, text=white, minimum size=2.5cm, font=\bfseries\large, drop shadow={opacity=0.2}},
  choix/.style={rectangle, rounded corners=5pt, draw=#1, line width=0.8pt, fill=white,
                text width=4.1cm, align=center, font=\small, inner sep=6pt, minimum height=1.55cm},
  lien/.style={draw=#1, line width=1pt, -{Latex[length=2.2mm]}},
  ]
  \node[centre] (e) at (0,0) {EquiPay};
  \node[choix=lagon]     (a) at (-5.3, 2.15) {\textbf{Mobile}\\ Flutter, Material 3,\\ carte native Apple Plans};
  \node[choix=violet]    (b) at ( 0.0, 3.25) {\textbf{Données}\\ Supabase : PostgreSQL,\\ RLS, stockage, temps réel};
  \node[choix=orange]    (c) at ( 5.3, 2.15) {\textbf{Service d'IA}\\ FastAPI en local,\\ inférence MLX 4 bits};
  \node[choix=bleuclair] (d) at (-5.3,-2.15) {\textbf{Modèles}\\ Qwen3-4B et Qwen3-VL-4B\\ affinés par LoRA};
  \node[choix=recevoir]  (f) at ( 0.0,-3.25) {\textbf{MLOps}\\ MLflow : suivi des essais,\\ registre et « champion »};
  \node[choix=ambre]     (g) at ( 5.3,-2.15) {\textbf{Itinéraires}\\ OSRM, algorithme d'ordre\\ et MapKit};
  \foreach \n/\c in {a/lagon, b/violet, c/orange, d/bleuclair, f/recevoir, g/ambre}
    \draw[lien=\c] (e) -- (\n);
\end{tikzpicture}
\caption{Synthèse des choix technologiques et fonctionnels retenus}
\label{fig:choix}
\end{figure}

Ces choix constituent le socle technique et fonctionnel détaillé dans le chapitre suivant,
consacré à la conception de l'application.

\subsection{Conclusion du chapitre}

Ce chapitre a posé le cadre du projet à travers l'analyse de l'existant, l'étude des
solutions concurrentes et la définition des concepts clés (modèles de langage et
vision-langage, LoRA, quantification, MLOps). L'étude de faisabilité et la spécification
des besoins fonctionnels et non fonctionnels ont guidé les choix retenus, qui servent de
base à la conception présentée dans le chapitre suivant.

% =====================================================================
\chapitre{Conception de l'application}
% =====================================================================

Ce chapitre traduit les besoins identifiés en une conception concrète. Il décrit
l'architecture globale de l'application, le modèle de données et ses règles d'accès, les
diagrammes UML des principaux scénarios, le cycle de vie commun à tous les modèles d'IA et,
pour finir, les choix d'interface.

\subsection{Architecture globale}

\equipay repose sur une architecture découplée en trois blocs : l'application mobile, la
plateforme Supabase qui détient les données partagées, et un backend d'IA qui sert les
modèles. Un quatrième bloc, MLflow, ne participe pas au fonctionnement quotidien mais
conserve l'historique des entraînements et désigne la version de chaque modèle que le
backend doit charger.

\subsubsection{Application mobile --- Flutter}

L'application est développée avec Flutter~\cite{flutter} et suit les principes de Material 3.
Elle regroupe les écrans, la logique de présentation et une couche de services (activités,
dépenses, programme, géocodage, IA). Chaque service est défini par une interface, avec une
implémentation réelle (Supabase ou HTTP) et une implémentation en mémoire utilisée par les
tests. Sur iPhone, la carte et la recherche de lieux s'appuient directement sur MapKit~\cite{mapkit}
grâce à une vue native écrite en Swift.

\subsubsection{Données --- Supabase}

Supabase~\cite{supabase} fournit quatre services utilisés par l'application :
l'authentification (e-mail et Google), une base PostgreSQL protégée par des règles RLS, le
stockage des photos et des reçus, et la diffusion en temps réel des modifications. L'application
lit et écrit directement dans la base : aucune donnée ne transite par le backend d'IA.

\subsubsection{Service d'IA --- FastAPI}

Le backend, écrit en Python avec FastAPI~\cite{fastapi}, expose quatre routes, une par
fonction d'IA. Il s'exécute sur un Mac Apple Silicon et fait tourner les modèles avec
MLX~\cite{mlx}. Chaque requête doit porter le jeton JWT de l'utilisateur, que le backend
vérifie avec la clé publique de Supabase. Les modèles sont chargés à la première
utilisation, puis gardés en mémoire.

\subsubsection{Entraînement et suivi --- MLflow}

Les modèles sont entraînés sur un PC équipé d'une RTX 5080, puis importés sur le Mac.
Le serveur MLflow~\cite{mlflow_docs}, installé sur le Mac et accessible sur le réseau local,
reçoit les paramètres, les métriques et les adaptateurs de chaque entraînement. Son registre
attribue l'alias « champion » à la meilleure version, que le backend charge au démarrage.

\subsubsection{Flux global}

La figure~\ref{fig:architecture} montre comment ces blocs communiquent.

\begin{figure}[H]
\centering
\resizebox{\textwidth}{!}{%
\begin{tikzpicture}[
  bloc/.style={rectangle, rounded corners=5pt, draw=#1, line width=1pt, fill=#1!7, align=center,
               text width=4.0cm, inner sep=7pt, font=\small},
  ext/.style={rectangle, rounded corners=5pt, draw=gristexte, dashed, fill=grisfond, align=center,
              text width=3.3cm, inner sep=6pt, font=\small},
  fl/.style={-{Latex[length=2.4mm]}, line width=0.9pt, draw=#1},
  et/.style={font=\scriptsize, fill=white, inner sep=1.5pt, align=center},
  ]
  \node[bloc=lagon] (app) at (0,-2.1) {\textbf{Application mobile}\\[2pt] Flutter (iPhone)\\[4pt]
     \scriptsize Écrans · services\\ \scriptsize carte native MapKit};
  \node[ext] (plans) at (0,-5.4) {\textbf{Apple Plans}\\ \scriptsize recherche de lieux, carte,\\ \scriptsize itinéraires};
  \node[bloc=violet] (sb) at (6.6,1.6) {\textbf{Supabase}\\[2pt] \scriptsize Auth · PostgreSQL + RLS\\
     \scriptsize Storage · Realtime};
  \node[bloc=orange] (api) at (6.6,-2.1) {\textbf{Backend d'IA}\\[2pt] FastAPI sur Mac M3 Pro\\[4pt]
     \scriptsize /chat · /receipts/scan\\ \scriptsize /expenses/describe · /trips/optimize\\[2pt]
     \scriptsize modèles MLX 4 bits};
  \node[ext] (osrm) at (12.6,-1.2) {\textbf{OSRM}\\ \scriptsize temps et distances\\ \scriptsize entre les lieux};
  \node[bloc=recevoir] (mlf) at (12.6,-3.4) {\textbf{MLflow}\\[2pt] \scriptsize essais, métriques,\\ \scriptsize registre « champion »};
  \node[bloc=bleuclair] (pc) at (12.6,-6.2) {\textbf{PC RTX 5080}\\[2pt] \scriptsize entraînement LoRA,\\ \scriptsize évaluation};

  \draw[fl=violet] (app.north) |- node[et, pos=0.75, above=1pt] {données, fichiers, temps réel} (sb.west);
  \draw[fl=orange] (app.east |- api.west) -- node[et, above=2pt] {HTTP + JWT} (api.west);
  \draw[fl=gristexte] (app.south) -- node[et, right=2pt] {MapKit} (plans.north);
  \draw[fl=violet, dashed] (api.north) -- node[et, right=2pt] {clé JWT, historique du chat} (sb.south);
  \draw[fl=gristexte] (api.east |- osrm.west) -- node[et, above=2pt] {temps de trajet} (osrm.west);
  \draw[fl=recevoir] (mlf.west) -- node[et, below=2pt] {version champion} (api.east |- mlf.west);
  \draw[fl=bleuclair] (pc.north) -- node[et, right=2pt] {runs, adaptateurs} (mlf.south);
\end{tikzpicture}}
\caption{Architecture globale et flux de données d'EquiPay}
\label{fig:architecture}
\end{figure}

Le parcours d'une donnée suit toujours la même logique : l'application envoie au backend
une photo, une phrase ou une liste de lieux ; le backend renvoie une proposition structurée
sans rien enregistrer ; l'utilisateur la vérifie, la corrige si besoin, et c'est
l'application qui l'enregistre dans Supabase. Les autres membres du groupe la reçoivent
alors en temps réel.

\begin{encadre}
\textbf{Règle de conception.} L'IA ne crée jamais rien à la place de l'utilisateur : le scanner
et la saisie par phrase pré-remplissent un formulaire, et le trajet proposé n'est enregistré
qu'après un appui sur « Appliquer ». La décision finale reste toujours humaine.
\end{encadre}

% ---------------------------------------------------------------------
\subsection{Conception de la base de données}

Le schéma a été conçu autour de l'activité, à laquelle se rattachent les membres, les
dépenses, les photos et le programme. Il est entièrement décrit par des fichiers de
migration SQL versionnés, appliqués dans l'ordre, ce qui permet de recréer la base à
l'identique. Le tableau~\ref{tab:tables} présente les tables principales.

\begin{table}[H]
\centering
\small
\begin{tabularx}{\textwidth}{|L{4.6cm}|Y|}
\hline
\rowcolor{entete}\textbf{Table} & \textbf{Contenu} \\ \hline
\code{profiles} & Profil de chaque utilisateur : nom d'utilisateur, nom complet, photo. \\ \hline
\code{activities} & Activité : type (voyage, sortie, quotidien, crédit), nom, devise, logo, dates. \\ \hline
\code{activity\_members} & Participants, rôle (créateur, membre, spectateur), date d'archivage personnelle. \\ \hline
\code{activity\_join\_codes}, \code{join\_requests} & Code d'accès à huit caractères et demandes pour rejoindre. \\ \hline
\code{expenses}, \code{expense\_shares} & Dépenses et part de chaque participant, exprimée en pourcentage. \\ \hline
\code{attachments} & Reçus et photos de l'activité ; les fichiers eux-mêmes sont dans le stockage. \\ \hline
\code{trip\_items}, \code{trip\_votes} & Programme du voyage et votes pour ou contre. \\ \hline
\code{chat\_messages} & Historique des échanges avec l'assistant. \\ \hline
\end{tabularx}
\caption{Tables principales de la base de données}
\label{tab:tables}
\end{table}

\subsubsection{Activités, membres et accès}

Chaque activité a ses propres rôles : le créateur gère tout, le membre ajoute des dépenses
et vote, le spectateur consulte. Une même personne peut être créatrice d'une activité et
simple membre d'une autre. L'archivage est personnel : il est stocké dans
\code{activity\_members} et ne masque l'activité que pour celui qui l'archive.

\subsubsection{Dépenses et répartitions}

Une dépense est enregistrée dans sa devise d'origine avec son taux de change vers la devise
de l'activité. La part de chaque participant est stockée en pourcentage, ce qui couvre tous
les modes de répartition (égale, montants, parts) avec une seule représentation. Une
fonction SQL vérifie que les parts d'une dépense totalisent exactement 100 \%, et une
fonction \code{save\_expense} enregistre la dépense et ses parts en une seule transaction.

\subsubsection{Programme du voyage}

La table \code{trip\_items} contient les lieux proposés, avec leur adresse, leur date,
leur heure et leurs coordonnées, remplies automatiquement à partir de l'adresse. Les votes
sont stockés dans \code{trip\_votes}, avec une seule ligne par personne et par lieu. Une
fonction \code{has\_program} réserve le programme aux activités de type voyage ou sortie.

\subsubsection{Sécurité au niveau des lignes}

Toutes les tables sont protégées par des règles RLS. Elles reposent sur quelques fonctions
réutilisées partout (\code{is\_member}, \code{is\_owner}, \code{can\_edit}) : un
utilisateur ne lit que les activités dont il est membre, ne modifie que ce qu'il a créé
(sauf le créateur de l'activité), et un spectateur ne peut rien modifier. Ces règles
s'appliquent quelle que soit la façon dont la base est interrogée.

% ---------------------------------------------------------------------
\subsection{Diagrammes UML}

Cette section présente les diagrammes UML qui modélisent l'application : le diagramme de
cas d'utilisation décrit les interactions entre les acteurs et le système, le diagramme de
classes représente la structure des données, et les diagrammes de séquence détaillent le
déroulement des quatre fonctions d'IA.

\subsubsection{Diagramme de cas d'utilisation}

\begin{figure}[H]
\centering
\resizebox{0.98\textwidth}{!}{%
\begin{tikzpicture}[
  uc/.style={ellipse, draw=lagon!80!black, fill=lagonclair!40, line width=0.7pt, minimum width=3.9cm,
             minimum height=0.8cm, font=\scriptsize, align=center, inner sep=1pt},
  ucx/.style={uc, draw=orange!80!black, fill=orangeclair!60},
  as/.style={line width=0.6pt, draw=black!65},
  dep/.style={dashed, -{Latex[length=2mm]}, draw=black!60, line width=0.6pt},
  st/.style={font=\tiny\itshape, fill=white, inner sep=1pt},
  ]
  % Acteur : bonhomme
  \newcommand{\acteur}[3]{%
    \begin{scope}[shift={(#1)}]
      \draw[line width=0.8pt] (0,0.55) circle (0.2);
      \draw[line width=0.8pt] (0,0.35) -- (0,-0.25) (-0.32,0.15) -- (0.32,0.15) (0,-0.25) -- (-0.25,-0.65) (0,-0.25) -- (0.25,-0.65);
      \node[font=\small\bfseries, align=center, anchor=north] at (0,-0.75) {#2};
      \coordinate (#3) at (0,0);
    \end{scope}}
  % Limite du système
  \draw[rounded corners=6pt, line width=0.9pt, draw=bleu, fill=grisfond!50] (1.25,0.35) rectangle (10.75,-10.55);
  \node[font=\bfseries\color{bleu}, anchor=north west] at (1.4,0.3) {EquiPay};

  % Cas du membre (colonne de gauche)
  \foreach \n/\t [count=\i] in {
      auth/{S'inscrire, se connecter},
      creer/{Créer une activité},
      rejoindre/{Rejoindre par code},
      depense/{Ajouter une dépense},
      soldes/{Consulter soldes et résumé},
      lieu/{Ajouter un lieu au programme},
      voter/{Voter pour ou contre},
      trajet/{Organiser le trajet},
      question/{Poser une question à l'assistant}}
    \node[uc] (\n) at (3.6,-0.9*\i-0.1) {\t};
  \node[uc] (accepter) at (3.6,-9.95) {Accepter une demande};

  % Extensions et inclusions (colonne de droite)
  \node[ucx] (scan) at (8.4,-2.75) {Scanner un ticket};
  \node[ucx] (phrase) at (8.4,-3.85) {Décrire la dépense\\ en une phrase};
  \node[ucx] (carte) at (8.4,-5.5) {Placer le lieu sur la carte};
  \node[ucx] (appliquer) at (8.4,-8.95) {Appliquer le programme};

  \draw[dep] (scan) -- node[st, above, sloped] {«extend»} (depense);
  \draw[dep] (phrase) -- node[st, below, sloped] {«extend»} (depense);
  \draw[dep] (lieu) -- node[st, above] {«include»} (carte);
  \draw[dep] (appliquer) -- node[st, below, sloped] {«extend»} (trajet);

  \acteur{-0.9,-3.9}{Membre}{membre}
  \acteur{-0.9,-8.6}{Créateur}{createur}
  \acteur{12.9,-5.0}{Modèles\\ d'IA}{ia}

  \draw[-{Triangle[open, length=3mm]}, line width=0.7pt] ([yshift=0.9cm]createur) -- ([yshift=-1.75cm]membre);
  \foreach \n in {auth, creer, rejoindre, depense, soldes, lieu, voter, trajet, question}
    \draw[as] ([xshift=0.35cm]membre) -- (\n.west);
  \draw[as] ([xshift=0.35cm]createur) -- (accepter.west);
  \draw[as] ([xshift=0.35cm, yshift=-0.1cm]createur) -- (appliquer.west);
  \foreach \n in {scan, phrase} \draw[as] ([xshift=-0.35cm]ia) -- (\n.east);
  \draw[as] ([xshift=-0.35cm]ia) -- (trajet.east);
  \draw[as] ([xshift=-0.35cm]ia) -- (question.east);
\end{tikzpicture}}
\caption{Diagramme de cas d'utilisation de l'application}
\label{fig:cas_utilisation}
\end{figure}

La figure~\ref{fig:cas_utilisation} présente les interactions entre les acteurs et le
système. Le membre gère son compte, crée ou rejoint des activités, ajoute des dépenses,
consulte les soldes, propose des lieux, vote, demande l'organisation du trajet et
interroge l'assistant. Le créateur est un membre qui dispose en plus de deux droits :
accepter les demandes pour rejoindre et appliquer le programme proposé. Les modèles d'IA
interviennent comme acteur secondaire ; la carte et le calcul des trajets s'appuient en
plus sur Apple Plans et OSRM, qui ne sont pas représentés pour alléger le schéma.

\subsubsection{Diagramme de classes}

\begin{figure}[H]
\centering
\resizebox{\textwidth}{!}{%
\begin{tikzpicture}[
  classe/.style={rectangle split, rectangle split parts=2, draw=bleu, line width=0.7pt,
                 rectangle split part fill={bleu!12, white}, text width=3.55cm, align=left,
                 font=\scriptsize, inner sep=3.5pt},
  as/.style={line width=0.7pt, draw=bleu!80},
  comp/.style={as, {Diamond[length=3.2mm, width=2.2mm]}-},
  mult/.style={font=\tiny, inner sep=1pt},
  ]
  \newcommand{\titre}[1]{\centering\textbf{\small #1}}
  \node[classe] (jr) at (0,0) {\titre{JoinRequest}\nodepart{two}
     id : uuid\\ activity\_name : text\\ status : pending | accepted\\ \hspace{2.6em}| declined\\ created\_at, decided\_at};
  \node[classe] (pr) at (5.4,0) {\titre{Profile}\nodepart{two}
     id : uuid\\ username : text\\ email : text\\ full\_name : text\\ avatar\_url : text};
  \node[classe] (ch) at (10.8,0) {\titre{ChatMessage}\nodepart{two}
     id : uuid\\ role : user | assistant\\ content : text\\ created\_at : timestamptz};

  \node[classe] (ex) at (0,-4.7) {\titre{Expense}\nodepart{two}
     id : uuid\\ name : text\\ amount : numeric\\ currency : char(3)\\ exchange\_rate : numeric\\ expense\_date : date\\ paid\_by : Profile};
  \node[classe] (ac) at (5.4,-4.7) {\titre{Activity}\nodepart{two}
     id : uuid\\ type : voyage | sortie |\\ \hspace{2.6em}quotidien | crédit\\ name : text\\ default\_currency : char(3)\\ start\_date, end\_date\\ cover\_url : text};
  \node[classe] (am) at (10.8,-4.7) {\titre{ActivityMember}\nodepart{two}
     role : owner | member\\ \hspace{2.6em}| viewer\\ joined\_at : timestamptz\\ archived\_at : timestamptz};

  \node[classe] (es) at (-0.6,-9.6) {\titre{ExpenseShare}\nodepart{two}
     user\_id : Profile\\ percentage : numeric};
  \node[classe] (at) at (3.4,-9.6) {\titre{Attachment}\nodepart{two}
     id : uuid\\ kind : receipt | photo\\ storage\_path : text};
  \node[classe] (tv) at (7.4,-9.6) {\titre{TripVote}\nodepart{two}
     user\_id : Profile\\ vote : -1 | 1};
  \node[classe] (ti) at (11.4,-9.6) {\titre{TripItem}\nodepart{two}
     id : uuid\\ name, place, address\\ planned\_on : date\\ planned\_at : time\\ latitude, longitude\\ done : boolean};

  \draw[as] (jr) -- node[mult, pos=0.12, above] {*} node[mult, pos=0.88, above] {1} (pr);
  \draw[as] (pr) -- node[mult, pos=0.12, above] {1} node[mult, pos=0.88, above] {*} (ch);
  \draw[as] (pr) -- node[mult, pos=0.15, right] {1} node[mult, pos=0.5, right, font=\tiny\itshape] {crée} node[mult, pos=0.85, right] {*} (ac);
  \draw[as] (jr.south east) -- node[mult, pos=0.1, right] {*} node[mult, pos=0.9, left] {1} (ac.north west);
  \draw[as] (pr.south east) -- node[mult, pos=0.1, right] {1} node[mult, pos=0.9, above] {*} (am.north west);
  \draw[comp] (ac) -- node[mult, pos=0.85, above] {*} (am);
  \draw[comp] (ac) -- node[mult, pos=0.85, above] {*} (ex);
  \draw[comp] (ex) -- node[mult, pos=0.85, left] {1..*} (es.north -| ex);
  \draw[comp] (ac.south west) -- node[mult, pos=0.85, right] {*} (at.north);
  \draw[as] (ex.south east) -- node[mult, pos=0.1, above] {0..1} node[mult, pos=0.85, left] {*} (at.north west);
  \draw[comp] (ac.south east) -- node[mult, pos=0.88, above] {*} (ti.north west);
  \draw[comp] (ti) -- node[mult, pos=0.8, above] {*} (tv);
\end{tikzpicture}}
\caption{Diagramme de classes du modèle de données}
\label{fig:classes}
\end{figure}

La figure~\ref{fig:classes} modélise la structure des données, directement dérivée du
schéma PostgreSQL. L'activité est au centre : elle se compose de ses membres, de ses
dépenses, de ses photos et de son programme, qui disparaissent avec elle. Chaque dépense
se décompose en parts, une par participant, et peut être liée à la photo de son reçu.
Les demandes pour rejoindre et l'historique du chat sont rattachés au profil de
l'utilisateur.

\subsubsection{Diagrammes de séquence}

Les diagrammes suivants détaillent le déroulement des quatre fonctions d'IA, de l'action
de l'utilisateur jusqu'à l'enregistrement du résultat.

\begin{figure}[H]
\centering
\resizebox{0.97\textwidth}{!}{\footnotesize%
\begin{sequencediagram}
  \gdef\unitfactor{0.82}
  \tikzstyle{inststyle}+=[fill=lagonclair!50, rounded corners=2pt]
  \newthread[lagonclair]{u}{:Utilisateur}
  \newinst[2.9]{app}{:Application}
  \newinst[3.3]{api}{:Backend FastAPI}
  \newinst[3.3]{vlm}{:Qwen3-VL-4B (MLX)}
  \newinst[2.5]{sb}{:Supabase}
  \begin{call}{u}{photographie le ticket}{app}{formulaire pré-rempli}
    \begin{call}{app}{POST /receipts/scan}{api}{champs de la dépense}
      \begin{callself}{api}{vérifie le jeton}{}
      \end{callself}
      \begin{call}{api}{lit l'image}{vlm}{JSON brut}
      \end{call}
      \begin{callself}{api}{normalise montant, devise, date}{}
      \end{callself}
    \end{call}
  \end{call}
  \begin{call}{u}{vérifie, corrige et valide}{app}{dépense ajoutée}
    \begin{call}{app}{save\_expense}{sb}{ok}
    \end{call}
  \end{call}
  \mess[1]{sb}{temps réel : nouvelle dépense}{app}
\end{sequencediagram}}
\caption{Diagramme de séquence --- Scanner un ticket de caisse}
\label{fig:seq_scan}
\end{figure}

La figure~\ref{fig:seq_scan} illustre la lecture d'un ticket : l'application envoie la photo
au backend, qui vérifie le jeton de l'utilisateur, fait lire l'image par le modèle
vision-langage puis normalise sa réponse. Le formulaire est pré-rempli, et la dépense n'est
enregistrée dans Supabase qu'après validation par l'utilisateur.

\begin{figure}[H]
\centering
\resizebox{0.97\textwidth}{!}{\footnotesize%
\begin{sequencediagram}
  \gdef\unitfactor{0.82}
  \tikzstyle{inststyle}+=[fill=lagonclair!50, rounded corners=2pt]
  \newthread[lagonclair]{u}{:Utilisateur}
  \newinst[2.9]{app}{:Application}
  \newinst[3.3]{api}{:Backend FastAPI}
  \newinst[3.3]{llm}{:equipay-describe (MLX)}
  \newinst[2.5]{sb}{:Supabase}
  \begin{call}{u}{écrit sa phrase}{app}{formulaire pré-rempli}
    \begin{call}{app}{POST /expenses/describe}{api}{dépense structurée}
      \begin{call}{api}{phrase, membres, date}{llm}{JSON brut}
      \end{call}
      \begin{callself}{api}{relie les noms aux membres}{}
      \end{callself}
      \begin{callself}{api}{parts ramenées à 100 \%}{}
      \end{callself}
    \end{call}
  \end{call}
  \begin{call}{u}{vérifie et valide}{app}{dépense ajoutée}
    \begin{call}{app}{save\_expense}{sb}{ok}
    \end{call}
  \end{call}
\end{sequencediagram}}
\caption{Diagramme de séquence --- Décrire une dépense en une phrase}
\label{fig:seq_phrase}
\end{figure}

La figure~\ref{fig:seq_phrase} décrit la saisie par phrase. Le modèle reçoit la phrase, la
liste des membres et la date du jour ; le backend relie ensuite chaque nom cité à un membre
réel, même mal orthographié, et convertit la répartition en pourcentages dont la somme fait
exactement 100.

\begin{figure}[H]
\centering
\resizebox{0.97\textwidth}{!}{\footnotesize%
\begin{sequencediagram}
  \gdef\unitfactor{0.82}
  \tikzstyle{inststyle}+=[fill=lagonclair!50, rounded corners=2pt]
  \newthread[lagonclair]{u}{:Utilisateur}
  \newinst[2.7]{app}{:Application}
  \newinst[3.1]{api}{:Backend FastAPI}
  \newinst[2.7]{osrm}{:OSRM}
  \newinst[2.7]{llm}{:Qwen3-4B (MLX)}
  \newinst[2.0]{sb}{:Supabase}
  \begin{call}{u}{choisit les options}{app}{trajet proposé}
    \begin{call}{app}{POST /trips/optimize}{api}{ordre, heures, gain, texte}
      \begin{call}{api}{matrice des temps entre lieux}{osrm}{durées, distances}
      \end{call}
      \begin{callself}{api}{répartit par jour, meilleur ordre}{}
      \end{callself}
      \begin{call}{api}{programme calculé à rédiger}{llm}{texte}
      \end{call}
    \end{call}
  \end{call}
  \begin{call}{u}{appliquer (créateur)}{app}{programme mis à jour}
    \begin{call}{app}{date et heure de chaque lieu}{sb}{ok}
    \end{call}
  \end{call}
\end{sequencediagram}}
\caption{Diagramme de séquence --- Organiser le trajet d'une journée}
\label{fig:seq_trajet}
\end{figure}

La figure~\ref{fig:seq_trajet} présente l'organisation du trajet. Les temps de parcours sont
calculés par OSRM et l'ordre de visite par un algorithme ; le modèle de langage ne fait que
rédiger le programme à partir de ces résultats, sans pouvoir modifier l'ordre ni les heures.
Seul le créateur de l'activité peut appliquer le programme proposé.

\begin{figure}[H]
\centering
\resizebox{0.86\textwidth}{!}{\footnotesize%
\begin{sequencediagram}
  \gdef\unitfactor{0.82}
  \tikzstyle{inststyle}+=[fill=lagonclair!50, rounded corners=2pt]
  \newthread[lagonclair]{u}{:Utilisateur}
  \newinst[2.9]{app}{:Application}
  \newinst[3.3]{api}{:Backend FastAPI}
  \newinst[3.3]{llm}{:equipay-assistant (MLX)}
  \newinst[2.5]{sb}{:Supabase}
  \begin{call}{u}{pose une question}{app}{réponse affichée}
    \begin{call}{app}{POST /chat (question, JWT)}{api}{réponse}
      \begin{call}{api}{question + prompt système}{llm}{réponse générée}
      \end{call}
      \begin{call}{api}{enregistre l'échange}{sb}{ok}
      \end{call}
    \end{call}
  \end{call}
\end{sequencediagram}}
\caption{Diagramme de séquence --- Question à l'assistant}
\label{fig:seq_chat}
\end{figure}

La figure~\ref{fig:seq_chat} montre l'échange avec l'assistant. Le backend génère la réponse
avec le modèle champion puis enregistre la question et la réponse dans Supabase, au nom de
l'utilisateur, afin de conserver l'historique de la conversation.

% ---------------------------------------------------------------------
\subsection{Architecture fonctionnelle des modules}

L'application s'articule autour de quatre modules fonctionnels, chacun responsable d'une
partie du parcours de l'utilisateur.

\subsubsection{Module des dépenses et des soldes}

Ce module gère les activités, les dépenses et leur répartition. Le solde de chaque
participant est la différence entre ce qu'il a payé et la somme de ses parts, convertie
dans la devise de l'activité. À partir de ces soldes, l'application calcule la liste des
remboursements qui remettent tous les comptes à zéro avec le moins de virements possible,
et propose des graphiques et des statistiques sur les dépenses du groupe.

\subsubsection{Module de saisie assistée}

Ce module regroupe les deux façons d'ajouter une dépense sans tout taper : le scanner de
tickets et la description en une phrase. Les deux aboutissent au même formulaire, où chaque
champ rempli par l'IA est signalé avec son origine (« Lu sur le ticket », « Compris dans ta
phrase ») et où les champs manquants sont mis en évidence.

\subsubsection{Module du programme et du trajet}

Ce module couvre la liste des lieux à visiter, les votes, l'affichage en liste, en
calendrier et sur une carte, ainsi que l'organisation du trajet. Les lieux sont placés
automatiquement sur la carte à partir de leur nom ou de leur adresse grâce à la recherche
d'Apple Plans.

\subsubsection{Module d'assistance}

L'assistant est accessible depuis une bulle « Support » présente sur les écrans principaux.
Il répond aux questions sur le fonctionnement de l'application, oriente vers le support
quand il ne sait pas, et refuse poliment les questions sans rapport avec \equipay.

% ---------------------------------------------------------------------
\subsection{Conception de la chaîne d'entraînement et de mise en service}

Les trois modèles affinés suivent le même cycle de vie, représenté dans la
figure~\ref{fig:mlops}. Ce cycle garantit qu'une nouvelle version n'est servie à
l'application que si elle fait mieux que la précédente sur le même jeu de test.

\begin{figure}[H]
\centering
\begin{tikzpicture}[
  etape/.style={rectangle, rounded corners=6pt, fill=#1, text=white, align=center, font=\small\bfseries,
                text width=3.3cm, minimum height=1.35cm, drop shadow={opacity=0.18}},
  det/.style={font=\scriptsize, align=center, text width=3.5cm, text=black!75},
  fl/.style={-{Latex[length=2.6mm]}, line width=1.1pt, draw=gristexte},
  ]
  \node[etape=lagon]     (d) at (0,2.2)     {1. Données\\ versionnées};
  \node[etape=violet]    (t) at (5.1,2.2)   {2. Entraînement\\ LoRA sur GPU};
  \node[etape=orange]    (e) at (10.2,2.2)  {3. Évaluation\\ comparative};
  \node[etape=recevoir]  (r) at (10.2,-1.6) {4. Registre MLflow\\ alias « champion »};
  \node[etape=bleuclair] (c) at (5.1,-1.6)  {5. Conversion\\ MLX 4 bits};
  \node[etape=ambre]     (s) at (0,-1.6)    {6. Service\\ FastAPI};
  \node[det, below=2pt of d] {génération, empreinte\\ du jeu enregistrée};
  \node[det, below=2pt of t] {PC RTX 5080, arrêt\\ anticipé sur la validation};
  \node[det, below=2pt of e] {base, ancienne et nouvelle\\ version sur le même test};
  \node[det, below=2pt of r] {seulement si la nouvelle\\ version est meilleure};
  \node[det, below=2pt of c] {fusion de l'adaptateur,\\ réévaluation sur le Mac};
  \node[det, below=2pt of s] {chargement au premier\\ appel, PyTorch en secours};
  \draw[fl] (d) -- (t);
  \draw[fl] (t) -- (e);
  \draw[fl] (e) -- (r);
  \draw[fl] (r) -- (c);
  \draw[fl] (c) -- (s);
  \draw[fl, dashed] (s.north) -- node[font=\scriptsize\itshape, fill=white, inner sep=2pt, align=center] {limites\\ observées} (d.south);
\end{tikzpicture}
\caption{Cycle de vie d'un modèle, des données à la mise en service}
\label{fig:mlops}
\end{figure}

\subsubsection{Construction des jeux de données}

Aucun jeu public ne correspondait à nos besoins pour l'assistant et la saisie par phrase.
Nous les avons donc construits à partir de sources écrites à la main (questions, réponses
de référence, phrases types) puis enrichis par génération de variantes. Pour le scanner, les
jeux publics de tickets ont été convertis dans notre format et complétés par des tickets
synthétiques dans les devises qui nous intéressent. Chaque jeu est découpé en
entraînement, validation et test, et son empreinte est enregistrée avec chaque essai.

\subsubsection{Entraînement}

Les modèles sont affinés par LoRA sur le PC, avec la bibliothèque PEFT~\cite{peft}. La
perte sur le jeu de validation est mesurée régulièrement, et l'entraînement s'arrête
lorsqu'elle cesse de s'améliorer. C'est l'adaptateur du meilleur point de validation qui est
conservé, et non le dernier.

\subsubsection{Évaluation et choix du champion}

Chaque nouvelle version est comparée au modèle de base et à la version précédente, sur le
même jeu de test et avec exactement le même code de calcul des métriques. Les métriques sont
regroupées en un score global pondéré, propre à chaque tâche, qui sert à décider. Si la
nouvelle version l'emporte, elle reçoit l'alias « champion » dans le registre MLflow.

\subsubsection{Conversion et mise en service}

L'adaptateur est fusionné avec le modèle de base, puis le modèle obtenu est converti en
4 bits avec MLX et réévalué sur le Mac, pour mesurer précisément ce que coûte la
quantification. Le backend charge ensuite la version champion. Pour le scanner, la version
PyTorch en pleine précision est conservée et sert de secours automatique si le chargement
MLX échoue.

% ---------------------------------------------------------------------
\subsection{Conception des interfaces utilisateur}

\subsubsection{Charte graphique}

L'identité visuelle, baptisée « Lagon », associe un sarcelle apaisant, qui évoque la
confiance et le voyage, à un orange énergique réservé aux actions principales. Les couleurs
des soldes sont les mêmes partout : vert pour ce que l'on doit recevoir, rouge pour ce que
l'on doit rembourser, gris lorsque l'on est à jour. Chaque type d'activité a sa propre
couleur (figure~\ref{fig:charte}). Le logo, un « e » sarcelle qui s'envole en avion au-dessus
d'un porte-monnaie orange, réunit le voyage et le partage des dépenses.

\begin{figure}[H]
\centering
\begin{tikzpicture}[
  pastille/.style={rectangle, rounded corners=4pt, minimum width=1.75cm, minimum height=1.05cm, fill=#1},
  nom/.style={font=\scriptsize\bfseries, align=center, text width=1.9cm},
  hex/.style={font=\tiny\ttfamily, text=gristexte},
  ]
  \foreach \c/\n/\h [count=\i] in {lagon/Principale/0D9488, orange/Accent/F97316, recevoir/À recevoir/16A34A,
                                  rembourser/À rembourser/DC2626, gristexte/À jour/64748B}
  {
    \node[pastille=\c] (p\i) at (2.15*\i,0) {};
    \node[nom, below=3pt of p\i] (n\i) {\n};
    \node[hex, below=-1pt of n\i] {\#\h};
  }
  \foreach \c/\n/\h [count=\i] in {lagon/Voyage/0D9488, violet/Sortie/8B5CF6, bleuclair/Quotidien/3B82F6, ambre/Crédit/F59E0B}
  {
    \node[pastille=\c, minimum width=1.5cm] (q\i) at (1.2+2.6*\i,-2.65) {};
    \node[nom, below=3pt of q\i] (m\i) {\n};
    \node[hex, below=-1pt of m\i] {\#\h};
  }
  \node[font=\small\bfseries\color{bleu}, anchor=east] at (1.0,0) {Couleurs};
  \node[font=\small\bfseries\color{bleu}, anchor=east, align=right] at (1.0,-2.65) {Types\\ d'activité};
\end{tikzpicture}
\caption{Palette « Lagon » : couleurs principales et couleurs des types d'activité}
\label{fig:charte}
\end{figure}

\subsubsection{Organisation des écrans}

La navigation repose sur une barre d'onglets en bas de l'écran : les activités, les demandes
pour rejoindre et le profil (un quatrième onglet est réservé à une évolution future). Depuis une activité, l'utilisateur accède à ses dépenses, à ses
membres, à ses photos et à son programme ; le menu à trois points regroupe les actions moins
fréquentes (modifier, archiver, résumé, graphiques, statistiques, supprimer). La
figure~\ref{fig:navigation} résume ce parcours.

\begin{figure}[H]
\centering
\resizebox{\textwidth}{!}{%
\begin{tikzpicture}[
  ecr/.style={rectangle, rounded corners=4pt, draw=#1, line width=0.8pt, fill=#1!8, font=\scriptsize,
              align=center, minimum height=0.75cm, text width=2.5cm},
  fl/.style={-{Latex[length=2mm]}, draw=gristexte, line width=0.7pt},
  ]
  \node[ecr=bleu] (tab) at (0,0) {\textbf{Barre d'onglets}};
  \node[ecr=lagon] (acc) at (-4.4,-1.5) {Accueil\\ (mes activités)};
  \node[ecr=lagon] (dem) at (0,-1.5) {Demandes\\ reçues / envoyées};
  \node[ecr=lagon] (pro) at (4.4,-1.5) {Profil\\ et préférences};
  \node[ecr=violet] (act) at (-4.4,-3.2) {Page de l'activité};
  \node[ecr=orange] (dep) at (-8.2,-4.9) {Dépenses\\ (ajout, détail)};
  \node[ecr=orange] (mem) at (-5.4,-4.9) {Membres et photos};
  \node[ecr=orange] (prg) at (-2.6,-4.9) {Programme};
  \node[ecr=orange] (men) at (0.4,-4.9) {Menu à trois points\\ résumé, graphiques};
  \node[ecr=ambre] (aj) at (-8.2,-6.6) {Saisie, scanner,\\ phrase};
  \node[ecr=ambre] (vues) at (-4.0,-6.6) {Liste, calendrier,\\ carte};
  \node[ecr=ambre] (traj) at (-1.2,-6.6) {Trajet proposé};
  \node[ecr=recevoir] (sup) at (4.4,-3.2) {Assistant\\ (bulle Support)};
  \foreach \a/\b in {tab/acc, tab/dem, tab/pro, acc/act, act/dep, act/mem, act/prg, act/men, dep/aj, prg/vues, prg/traj, pro/sup}
    \draw[fl] (\a) -- (\b);
\end{tikzpicture}}
\caption{Organisation des principaux écrans de l'application}
\label{fig:navigation}
\end{figure}

\subsubsection{Expérience utilisateur}

L'expérience a été pensée pour un usage rapide, souvent debout ou en déplacement. Les
actions principales sont toujours au même endroit (bouton orange en bas à droite), chaque
action donne un retour immédiat (chargement, confirmation, erreur expliquée), et les écrans
vides indiquent ce qu'il faut faire pour commencer. L'application propose un mode sombre,
fonctionne hors ligne en affichant un bandeau, et se met à jour d'elle-même lorsqu'un autre
membre ajoute une dépense ou vote.

% ---------------------------------------------------------------------
\subsection{Conclusion du chapitre}

Ce chapitre a posé les bases de la conception d'\equipay : une architecture qui sépare
clairement l'application, les données et l'IA, un modèle de données protégé par des règles
d'accès strictes, des diagrammes UML qui formalisent les scénarios clés, et un cycle de vie
commun à tous les modèles. Cette conception sert de fondation à la réalisation présentée
dans le chapitre suivant.

% =====================================================================
\chapitre{Réalisation et implémentation}
% =====================================================================

Ce chapitre présente la mise en œuvre de la conception : l'environnement de travail, le
choix des modèles, la construction, l'entraînement et l'évaluation de chacun d'eux,
l'organisation du trajet, puis les fonctionnalités transverses et la stratégie
d'exécution retenue.

Avant de détailler chaque fonctionnalité, il faut souligner une partie du travail qui ne se
voit pas dans l'application : l'outillage qui encadre le développement. Un script lance
l'ensemble des serveurs dans le bon ordre et un autre les arrête ; un autre enchaîne la
préparation des données, l'entraînement et l'évaluation d'un modèle ; la configuration
(adresse du backend, clés publiques) est centralisée dans un fichier unique passé à la
compilation de l'application ; les changements de base de données passent toujours par une
nouvelle migration ; et chaque modification est vérifiée par une suite de tests
automatiques. Ces éléments ne sont pas des fonctionnalités, mais ce sont eux qui nous ont
permis d'avancer vite sans casser ce qui fonctionnait déjà.

% ---------------------------------------------------------------------
\subsection{Environnement de travail et outils}

\subsubsection{Stack technique}

Le tableau~\ref{tab:stack} récapitule les technologies utilisées dans chaque partie du
projet.

\begin{table}[H]
\centering
\small
\begin{tabularx}{\textwidth}{|L{2.6cm}|L{4.6cm}|Y|}
\hline
\rowcolor{entete}\textbf{Couche} & \textbf{Technologie} & \textbf{Rôle} \\ \hline
Mobile & Flutter (Dart), Material 3 & Un seul code pour iOS et Android, interface moderne. \\ \hline
Carte & MapKit (vue native Swift) & Carte Apple Plans, recherche de lieux, itinéraires. \\ \hline
Données & Supabase : PostgreSQL, Auth, Storage, Realtime & Comptes, données partagées, fichiers et
mises à jour en direct. \\ \hline
Service d'IA & FastAPI (Python) & Routes typées, documentation générée automatiquement. \\ \hline
Inférence & MLX, \code{mlx-lm}, \code{mlx-vlm} & Exécution des modèles en 4 bits sur la puce du Mac. \\ \hline
Entraînement & PyTorch~\cite{pytorch}, Transformers, PEFT & Affinage LoRA sur la RTX 5080. \\ \hline
MLOps & MLflow & Suivi des essais, registre des modèles, alias « champion ». \\ \hline
Itinéraires & OSRM & Temps et distances réels entre les lieux. \\ \hline
\end{tabularx}
\caption{Stack technique du projet}
\label{tab:stack}
\end{table}

Dans l'application comme dans le backend, le code est organisé en couches qui ne dépendent
que de la couche inférieure (figure~\ref{fig:couches}). Côté mobile, les écrans ne parlent
qu'à des services, et chaque service a deux implémentations : la vraie (Supabase, HTTP,
géocodage natif) et une version en mémoire. Les tests utilisent la seconde, ce qui permet de
vérifier tous les écrans sans réseau. Côté backend, chaque route délègue à un module métier
qui charge et interroge son modèle.

\begin{figure}[H]
\centering
\begin{tikzpicture}[
  couche/.style={rectangle, rounded corners=4pt, fill=#1, text=white, font=\small\bfseries,
                 minimum width=6.4cm, minimum height=0.85cm, align=center},
  note/.style={font=\scriptsize, text=gristexte, align=left, anchor=west},
  fl/.style={-{Latex[length=2mm]}, line width=0.8pt, draw=gristexte},
  ]
  \node[font=\small\bfseries\color{bleu}] at (0,0.85) {Application mobile};
  \node[couche=lagon] (m1) at (0,0) {Écrans (pages, widgets)};
  \node[couche=lagon!80!black] (m2) at (0,-1.25) {Services (interfaces)};
  \node[couche=lagon!60!black] (m3) at (0,-2.5) {Supabase · HTTP · natif iOS};
  \node[couche=gristexte] (m4) at (0,-3.75) {Implémentations en mémoire (tests)};
  \draw[fl] (m1) -- (m2); \draw[fl] (m2) -- (m3);
  \draw[fl, dashed] (m2.south east) ++(-0.4,0) to[bend left=35] (m4.north east);

  \node[font=\small\bfseries\color{bleu}] at (7.6,0.85) {Backend d'IA};
  \node[couche=orange] (b1) at (7.6,0) {Routes FastAPI + vérification JWT};
  \node[couche=orange!80!black] (b2) at (7.6,-1.25) {Modules : assistant, receipts,\\ describe, trips};
  \node[couche=orange!60!black] (b3) at (7.6,-2.5) {Modèles MLX (chargement unique)};
  \node[couche=recevoir!80!black] (b4) at (7.6,-3.75) {Registre MLflow (champion)};
  \draw[fl] (b1) -- (b2); \draw[fl] (b2) -- (b3); \draw[fl] (b4) -- (b3);
\end{tikzpicture}
\caption{Organisation en couches de l'application et du backend}
\label{fig:couches}
\end{figure}

\subsubsection{Base de données et migrations}

Le schéma Supabase évolue uniquement par fichiers de migration datés, appliqués par un
script qui garde la trace de ceux déjà passés. Douze migrations ont été écrites au fil du
projet : le schéma initial, puis les photos, la suppression de compte, les demandes pour
rejoindre, le verrouillage de la devise, l'enregistrement atomique des dépenses,
l'archivage et enfin le programme des voyages. Une règle a été tenue du début à la fin : les
migrations ne contiennent que de la structure (tables, règles, fonctions), jamais de
données. La base livrée est vide, et ce sont les utilisateurs qui la remplissent.

\subsubsection{MLX et MLflow}

Le serveur MLflow tourne sur le Mac et sert à la fois l'interface de suivi et le stockage
des fichiers produits par les entraînements. Le PC d'entraînement y envoie ses résultats
par le réseau local. Chaque modèle a sa propre expérience (« Equipay » pour l'assistant,
« Equipay-receipts » pour le scanner, « Equipay-describe » pour la saisie par phrase,
« Equipay-trips » pour l'organisation des trajets) et son entrée dans le registre. Côté Mac,
les modèles sont convertis et exécutés avec MLX~\cite{mlx}, en utilisant les bibliothèques
\code{mlx-lm}~\cite{mlxlm} pour le texte et \code{mlx-vlm}~\cite{mlxvlm} pour la vision.

% ---------------------------------------------------------------------
\subsection{Choix des modèles d'intelligence artificielle}

Le projet utilise deux familles de modèles : un modèle de langage pour tout ce qui est
texte, et un modèle vision-langage pour les tickets. Les deux choix ont été guidés par la
même contrainte : un modèle assez petit pour être affiné sur une carte de 16 Go et servi
en local, mais assez bon en français pour comprendre des phrases familières.

\subsubsection{Modèle de langage : Qwen3-4B-Instruct-2507}

Pour l'assistant, nous avons d'abord comparé plusieurs modèles ouverts de petite taille,
tous affinés avec le même jeu de données (première version) et évalués avec le même code
(tableau~\ref{tab:choix_llm}). Sans affinage, aucun ne dépassait un score de 0,30 : ils
répondaient de façon générique et inventaient des fonctionnalités. Après affinage, Qwen3-4B
s'est montré le meilleur, et la version entraînée sur les 36 couches du modèle l'a encore
amélioré. C'est ce modèle qui a été retenu pour toutes les tâches de texte.

\begin{table}[H]
\centering
\small
\begin{tabularx}{\textwidth}{|L{4.4cm}|C{1.9cm}|C{2.2cm}|Y|C{0.9cm}|}
\hline
\rowcolor{entete}\textbf{Modèle} & \textbf{Paramètres} & \textbf{Score après affinage} & \textbf{Remarques} & \textbf{Réf.} \\ \hline
Qwen3-4B-Instruct-2507, LoRA sur les 36 couches & 4 milliards & \textbf{0,822} & Modèle retenu : meilleur score, bon
niveau en français. & \cite{qwen3_4b} \\ \hline
Qwen3-4B-Instruct-2507, LoRA partiel (Mac) & 4 milliards & 0,803 & Premier essai, LoRA sur une partie des couches, entraîné sur le Mac avec MLX. & \cite{qwen3_4b} \\ \hline
Qwen2.5-1.5B-Instruct & 1,5 milliard & 0,780 & Deux fois plus rapide, mais score inférieur. & \cite{qwen25} \\ \hline
Llama-3.2-3B-Instruct & 3 milliards & 0,770 & Score le plus faible des modèles affinés. & \cite{llama32} \\ \hline
\end{tabularx}
\caption{Comparaison des modèles de langage affinés pour l'assistant (jeu de données v1)}
\label{tab:choix_llm}
\end{table}

\subsubsection{Modèle vision-langage : Qwen3-VL-4B-Instruct}

Pour les tickets, nous avons écarté l'approche classique « OCR puis analyse du texte »,
qui perd la mise en page du ticket : la position d'un montant (en bas, en gras, à côté du
mot \textit{Total}) est justement ce qui permet de le reconnaître. Qwen3-VL-4B~\cite{qwen3vl}
lit directement l'image et répond en JSON. Sans affinage, il obtenait déjà un score de 0,792
sur notre jeu de test, ce qui en faisait un bon point de départ ; sa taille (4,5 milliards de
paramètres) reste compatible avec un affinage LoRA sur notre GPU.

\subsubsection{Inférence locale en 4 bits}

Avec le scanner servi en précision complète par PyTorch, le backend occupait environ 11 Go
de mémoire, ce qui laissait trop peu de place au reste du système sur un Mac de 18 Go. Nous avons donc converti chaque modèle en 4 bits avec MLX, puis mesuré l'effet de
cette conversion sur le même jeu de test (tableau~\ref{tab:quantification}). La perte est
faible et le gain de mémoire important : c'est la version 4 bits qui est servie, la version
PyTorch restant disponible en secours.

\begin{table}[H]
\centering
\small
\begin{tabularx}{\textwidth}{|Y|C{3.3cm}|C{3.3cm}|}
\hline
\rowcolor{entete}\textbf{Mesure (scanner de tickets)} & \textbf{PyTorch, 16 bits} & \textbf{MLX, 4 bits} \\ \hline
Score global & 0,954 & 0,941 \\ \hline
Total du ticket exact & 97,2 \% & 96,3 \% \\ \hline
Taille du modèle sur disque & 8,3 Go & 2,9 Go \\ \hline
Mémoire du backend (assistant et scanner chargés) & environ 11 Go & environ 5,8 Go \\ \hline
\end{tabularx}
\caption{Effet de la quantification en 4 bits sur le scanner de tickets}
\label{tab:quantification}
\end{table}

% ---------------------------------------------------------------------
\subsection{Assistant conversationnel}

\subsubsection{Jeu de données}

L'assistant doit adopter trois comportements : répondre aux questions sur l'application,
orienter vers le support quand la réponse n'existe pas (fonction absente, problème de
compte), et refuser poliment ce qui n'a rien à voir avec \equipay. Nous avons écrit à la main
les sources du jeu de données : 196 sujets répartis en douze catégories (activités,
dépenses, soldes, compte, thème, applications concurrentes, situations concrètes, etc.),
chacun avec une réponse de référence et des mots-clés attendus. Ces sources ont ensuite été
déclinées en 2\,423 questions formulées de manières différentes. Le jeu final compte
10\,921 exemples d'entraînement, 438 de validation et 727 de test.

Les réponses de référence décrivent l'état réel de l'application : lorsqu'une fonction n'est
pas encore développée (réinitialisation du mot de passe, invitation par e-mail, etc.), la
réponse attendue oriente vers le support plutôt que d'inventer une procédure.

\subsubsection{Entraînement}

Le modèle Qwen3-4B-Instruct-2507 a été affiné par LoRA sur ses 36 couches (rang 16,
$\alpha = 320$), sur le PC, pour cinq époques au maximum avec arrêt anticipé. La perte de
validation a atteint son minimum autour de la deuxième époque (0,052) : c'est ce point qui a
été conservé. La version obtenue est la seizième du registre \code{equipay-assistant}.

\subsubsection{Évaluation}

Une réponse générée est rarement identique mot pour mot à la référence. Nous avons donc
combiné plusieurs mesures complémentaires dans un score global :

\begin{itemize}
  \item le rappel des mots-clés attendus (poids 0,30), qui vérifie que les informations
        importantes sont présentes ;
  \item la similarité sémantique (poids 0,25), mesurée par le cosinus entre les plongements
        de la réponse et de la référence, calculés avec un modèle Sentence-BERT
        multilingue~\cite{sbert} ;
  \item le ROUGE-L~\cite{rouge} (poids 0,15), qui mesure la plus longue sous-séquence commune
        avec la référence ;
  \item le F1 macro du comportement (poids 0,30) : la réponse appartient-elle à la bonne
        classe (répondre, orienter vers le support, refuser) ?
\end{itemize}

Le BLEU~\cite{bleu} et le chrF~\cite{chrf} sont également calculés et enregistrés dans MLflow
à titre de comparaison, mais n'entrent pas dans le score. Le tableau~\ref{tab:assistant}
compare le modèle de base et la version retenue sur le jeu de test.

\begin{table}[H]
\centering
\small
\begin{tabularx}{\textwidth}{|Y|C{3.2cm}|C{3.2cm}|}
\hline
\rowcolor{entete}\textbf{Métrique (jeu de test, 727 questions)} & \textbf{Qwen3-4B de base} & \textbf{Champion v16} \\ \hline
Score global & 0,276 & \textbf{0,770} \\ \hline
Similarité sémantique & 0,496 & 0,807 \\ \hline
ROUGE-L & 0,146 & 0,697 \\ \hline
Comportement correct (exactitude) & 41,7 \% & 91,7 \% \\ \hline
Questions hors sujet auxquelles il répond quand même & 97,4 \% & 14,1 \% \\ \hline
Questions de support correctement orientées & --- & 70,5 \% \\ \hline
Temps de réponse moyen sur le Mac & --- & 1,3 s \\ \hline
\end{tabularx}
\caption{Résultats de l'assistant conversationnel avant et après affinage}
\label{tab:assistant}
\end{table}

Le gain est net sur toutes les mesures. Le modèle de base répondait à presque toutes les
questions, y compris aux recettes de cuisine, et décrivait des écrans qui n'existent pas.
Après affinage, il connaît l'application et respecte son rôle. La faiblesse restante concerne
les questions de support : sur une fonction inconnue, le modèle invente encore une réponse
dans environ 30 \% des cas au lieu d'orienter vers le support. C'est la piste d'amélioration
prioritaire pour une prochaine version du jeu de données.

% ---------------------------------------------------------------------
\subsection{Scanner de tickets de caisse}

\subsubsection{Une première version instructive}

La première version du scanner a été affinée uniquement sur CORD~\cite{cord}, un jeu de
tickets indonésiens. Son score sur ce jeu était excellent, mais en conditions réelles elle a
révélé deux défauts : elle ignorait les séparateurs décimaux (un total de 11,50 € devenait
1\,150) et annonçait systématiquement des roupies indonésiennes. Le modèle avait simplement
appris les habitudes de son unique source, où les montants sont des entiers et la devise
toujours la même. Cette erreur a orienté toute la seconde version.

\subsubsection{Jeu de données de la seconde version}

Le jeu de la seconde version réunit quatre sources : CORD, SROIE~\cite{sroie} (tickets
malaisiens), WildReceipt~\cite{wildreceipt} (photos prises dans des conditions variées) et
des tickets synthétiques que nous avons générés pour la France, le Maroc, les États-Unis, le
Royaume-Uni et la Suisse, avec des montants à virgule et des devises variées. Il compte
4\,638 tickets d'entraînement, 378 de validation et 615 de test, dans sept devises. Le modèle
doit renvoyer l'enseigne, l'adresse, la date, le total tel qu'il est écrit, le total en
nombre, la devise et la liste des articles. Comme toutes les sources n'annotent pas tous les
champs, les champs absents d'une source sont exclus du calcul de la perte, pour ne pas
apprendre au modèle à les laisser vides.

\subsubsection{Entraînement}

L'affinage LoRA (rang 16, $\alpha = 32$) porte sur les couches du modèle de langage
(projections d'attention et couches denses), la partie vision restant figée : 33 millions de
paramètres entraînés sur 4,47 milliards, soit 0,74 \%. Les images sont redimensionnées à
1\,024 pixels de côté au maximum, avec un lot effectif de 8 et trois époques au plus.
L'entraînement a été interrompu par une défaillance du GPU à l'étape 1\,595 sur 1\,740 ; le
meilleur point de validation ayant été atteint à l'étape 1\,158 et sauvegardé, aucune
perte de résultat n'en a découlé.

\subsubsection{Résultats}

Les trois modèles (base, v1 et v2) ont été évalués sur le même jeu de test
(tableau~\ref{tab:scanner} et figure~\ref{fig:scores}). Le score global pondère le total
(0,35), les articles (0,15), la devise, la date, l'enseigne et la validité du JSON (0,10
chacun), le total recopié et l'adresse (0,05 chacun).

\begin{table}[H]
\centering
\small
\begin{tabularx}{\textwidth}{|Y|C{2.4cm}|C{2.4cm}|C{2.4cm}|}
\hline
\rowcolor{entete}\textbf{Métrique (615 tickets de test)} & \textbf{Base} & \textbf{v1 (CORD seul)} & \textbf{v2 (retenue)} \\ \hline
Score global & 0,792 & 0,219 & \textbf{0,954} \\ \hline
Total exact & 71,2 \% & 21,1 \% & \textbf{97,2 \%} \\ \hline
Devise & 69,7 \% & 18,6 \% & \textbf{99,4 \%} \\ \hline
Date & 86,4 \% & 0,0 \% & \textbf{96,4 \%} \\ \hline
Enseigne & 89,4 \% & 0,0 \% & \textbf{91,3 \%} \\ \hline
Articles (F1) & 73,5 \% & 16,9 \% & \textbf{88,3 \%} \\ \hline
Ticket entièrement juste & 23,6 \% & 0,0 \% & \textbf{74,5 \%} \\ \hline
Erreur « virgule ignorée » & 0,0 \% & 72,7 \% & \textbf{0,0 \%} \\ \hline
\end{tabularx}
\caption{Comparaison du modèle de base, de la v1 et de la v2 du scanner}
\label{tab:scanner}
\end{table}

La v2 améliore toutes les mesures. Le total est exact sur 100 \% des tickets en euros et en
dirhams du jeu de test. L'analyse des erreurs du modèle de base est instructive : dans
13,5 \% des cas, il prenait pour total la somme donnée en espèces par le client, une
confusion que l'affinage a fait disparaître. Deux limites demeurent : le modèle invente
parfois une date lorsqu'aucune n'est imprimée, et environ 13 \% des articles proposés
n'existent pas sur le ticket. C'est pourquoi l'application indique l'origine de chaque champ
et laisse l'utilisateur corriger avant d'enregistrer.

% ---------------------------------------------------------------------
\subsection{Saisie d'une dépense en une phrase}

\subsubsection{Jeu de données et entraînement}

Le modèle reçoit une phrase libre (« Resto avec Sara, j'ai payé 45 € hier, 60/40 »), la
liste des membres de l'activité, la date du jour et la devise par défaut, et doit renvoyer
la dépense au format JSON : montant, devise, date, payeur, participants, répartition,
titre, enseigne et articles. Le jeu de données a été généré à partir de modèles de phrases
variés, dans quatre styles d'écriture (français, langage SMS, darija et anglais), avec
des prénoms écrits en minuscules ou avec de petites fautes : 12\,000 exemples
d'entraînement, 500 de validation et 800 de test, auxquels s'ajoutent 50 phrases écrites à
la main pour un test plus proche de l'usage réel. Le modèle est un Qwen3-4B-Instruct-2507
affiné par LoRA sur ses 36 couches pendant trois époques.

\subsubsection{Post-traitement}

Le modèle ne connaît que les noms écrits dans la phrase, alors que l'application a besoin
d'identifiants de membres. Le backend relie donc chaque nom à un membre réel, en tolérant
les fautes de frappe (figure~\ref{fig:code_matcher}) : « moi » désigne l'utilisateur
connecté, un nom exact est accepté directement, et sinon le membre le plus proche est
retenu seulement s'il est assez proche et sans ambiguïté. Un nom inconnu n'est jamais
inventé. La répartition est ensuite convertie en pourcentages dont la somme fait exactement
100 (figure~\ref{fig:code_pourcentages}).

\begin{figure}[H]
\begin{lstlisting}[language=Python]
def __call__(self, name) -> str | None:
    key = norm(name)
    if not key:
        return None
    if key in ("moi", "me", "ana", "je"):
        return self.me_id
    exact = {mid for k, mid in self.keys if k == key}
    if len(exact) == 1:
        return exact.pop()
    # Faute de frappe (« Slma » → Salma) : le nom le plus proche, s'il est assez proche et sans ambiguïté
    scored = sorted(((difflib.SequenceMatcher(None, key, k).ratio(), mid) for k, mid in self.keys), reverse=True)
    if scored and scored[0][0] >= 0.75 and all(mid == scored[0][1] for r, mid in scored if r >= scored[0][0] - 0.05):
        return scored[0][1]
    return None
\end{lstlisting}
\caption{Rattachement d'un nom cité dans la phrase à un membre de l'activité}
\label{fig:code_matcher}
\end{figure}

\begin{figure}[H]
\begin{lstlisting}[language=Python]
def to_percentages(weights: dict[str, float]) -> dict[str, float] | None:
    """Poids (parts, montants, %) → pourcentages au centième dont la somme fait exactement 100."""
    weights = {k: v for k, v in weights.items() if v > 0}
    total = sum(weights.values())
    if not weights or total <= 0:
        return None
    ids = list(weights)
    cents = [round(weights[i] / total * 10000) for i in ids]
    cents[-1] += 10000 - sum(cents)
    return {i: c / 100 for i, c in zip(ids, cents)}
\end{lstlisting}
\caption{Conversion d'une répartition quelconque en pourcentages exacts}
\label{fig:code_pourcentages}
\end{figure}

\subsubsection{Résultats}

\begin{table}[H]
\centering
\small
\begin{tabularx}{\textwidth}{|Y|C{2.7cm}|C{2.7cm}|C{2.7cm}|}
\hline
\rowcolor{entete}\textbf{Métrique (version MLX 4 bits)} & \textbf{Test (800)} & \textbf{Phrases manuelles (50)} & \textbf{Modèle de base} \\ \hline
Score global & \textbf{0,976} & \textbf{0,954} & 0,656 / 0,585 \\ \hline
Dépense entièrement juste & 72,5 \% & 66 \% & --- \\ \hline
Montant & 99,3 \% & 100 \% & --- \\ \hline
Participants (F1) & 0,994 & 0,959 & --- \\ \hline
Payeur & 97,2 \% & 97,8 \% & --- \\ \hline
Répartition & 98,8 \% & 97,8 \% & --- \\ \hline
Champs inventés & 1,6 \% & 15,6 \% & --- \\ \hline
\end{tabularx}
\caption{Résultats du modèle de saisie par phrase}
\label{tab:phrase}
\end{table}

Le tableau~\ref{tab:phrase} montre que le montant, le payeur et la répartition sont presque
toujours justes. Les phrases écrites à la main, plus libres que celles du jeu généré, font
apparaître davantage de champs inventés (souvent une enseigne ou un titre déduits à tort).
Deux cas restent difficiles : « pour nous deux » lorsqu'une autre personne a payé, et les
prix unitaires (« 2 cafés à 3,50 ») parfois compris comme le prix de la ligne. Sur le Mac, une
phrase est traitée en 1,5 à 2,5 secondes.

\begin{figure}[H]
\centering
\begin{tikzpicture}
\begin{axis}[
  ybar, bar width=0.62cm, width=0.86\textwidth, height=6.2cm,
  ymin=0, ymax=1.18, ytick={0,0.2,0.4,0.6,0.8,1},
  yticklabel style={/pgf/number format/.cd, fixed, fixed zerofill, precision=1, use comma},
  symbolic x coords={Assistant, Scanner de tickets, Saisie par phrase},
  xtick=data, x tick label style={font=\small},
  enlarge x limits=0.25,
  legend style={at={(0.5,1.04)}, anchor=south, legend columns=-1, draw=none, font=\small},
  nodes near coords, nodes near coords style={font=\scriptsize, /pgf/number format/.cd, fixed, fixed zerofill, precision=3, use comma},
  axis line style={draw=black!40}, tick style={draw=black!40},
  ymajorgrids, grid style={draw=black!8},
  ylabel={Score global}, ylabel style={font=\small},
  ]
  \addplot[fill=gristexte!60, draw=none] coordinates {(Assistant,0.276) (Scanner de tickets,0.792) (Saisie par phrase,0.656)};
  \addplot[fill=lagon, draw=none] coordinates {(Assistant,0.770) (Scanner de tickets,0.954) (Saisie par phrase,0.976)};
  \legend{Modèle de base, Modèle affiné (champion)}
\end{axis}
\end{tikzpicture}
\caption{Score global de chaque modèle avant et après affinage, sur son jeu de test}
\label{fig:scores}
\end{figure}

% ---------------------------------------------------------------------
\subsection{Organisation du trajet}

\subsubsection{Une chaîne en quatre étapes}

Confier entièrement l'organisation d'une journée à un modèle de langage aurait été risqué :
ces modèles calculent mal les distances et peuvent inventer des temps de trajet. Nous avons
donc séparé le calcul et la rédaction. Le modèle n'intervient qu'à la fin, pour expliquer un
résultat déjà calculé :

\begin{enumerate}
  \item \textbf{Temps de trajet.} Le backend demande à OSRM~\cite{osrm}, un moteur
        d'itinéraires fondé sur les données OpenStreetMap~\cite{osm}, la matrice des temps et
        des distances entre tous les lieux, à pied, à vélo ou en voiture. Si le service ne
        répond pas, les temps sont estimés à partir de la distance à vol d'oiseau, et
        l'application le signale.
  \item \textbf{Répartition par jour.} Les lieux datés restent sur leur jour. Les lieux sans
        date sont placés, du plus apprécié au moins apprécié, sur le jour dont ils sont le plus
        proches, dans la limite du nombre de visites choisi. Les lieux qui ont plus de votes
        contre que de votes pour sont écartés et listés à part.
  \item \textbf{Ordre de visite.} Pour chaque jour, le backend cherche l'ordre qui minimise le
        temps de déplacement, en respectant les heures déjà fixées par le groupe. Jusqu'à huit
        lieux, toutes les permutations sont testées ; au-delà, une heuristique du plus proche
        voisin est améliorée par la méthode 2-opt~\cite{croes} (figure~\ref{fig:code_ordre}).
        Des heures sont ensuite proposées pour les lieux qui n'en ont pas : départ à 9 h,
        90 minutes par visite, arrondi au quart d'heure.
  \item \textbf{Rédaction.} Qwen3-4B-Instruct-2507, sans affinage, reçoit le programme
        calculé et rédige un court texte par jour. Son prompt système lui interdit de
        modifier l'ordre, les heures ou les durées, et d'inventer un lieu ou une distance.
\end{enumerate}

\begin{figure}[H]
\begin{lstlisting}[language=Python]
def best_order(indices: list[int], stops: list[Stop], d: list[list[float]]) -> list[int]:
    """Ordre de visite d'une journée qui minimise le temps de trajet (chemin ouvert)."""
    if len(indices) <= EXACT_LIMIT:
        best, best_cost = None, math.inf
        for perm in itertools.permutations(indices):
            if not respects_times(list(perm), stops):
                continue
            cost = path_cost(list(perm), d)
            if cost < best_cost - 1e-9:
                best, best_cost = list(perm), cost
        return best or list(indices)
    # Heuristique : plus proche voisin depuis chaque départ, puis améliorations 2-opt
    best, best_cost = list(indices), math.inf
    for start in indices:
        order, rest = [start], set(indices) - {start}
        while rest:
            nxt = min(rest, key=lambda j: d[order[-1]][j])
            order.append(nxt)
            rest.remove(nxt)
        improved = True
        while improved:
            improved = False
            for i in range(1, len(order) - 1):
                for k in range(i + 1, len(order)):
                    candidate = order[:i] + order[i : k + 1][::-1] + order[k + 1 :]
                    if path_cost(candidate, d) < path_cost(order, d) - 1e-9 and respects_times(candidate, stops):
                        order, improved = candidate, True
        if respects_times(order, stops) and path_cost(order, d) < best_cost:
            best, best_cost = order, path_cost(order, d)
    return best
\end{lstlisting}
\caption{Recherche du meilleur ordre de visite d'une journée (extrait simplifié)}
\label{fig:code_ordre}
\end{figure}

\subsubsection{Garanties et suivi}

L'ordre proposé n'est jamais plus long que l'ordre actuel du groupe : s'il n'y a rien à
gagner, l'application l'indique (« Votre ordre est déjà le meilleur »). Chaque organisation
demandée est enregistrée dans l'expérience MLflow « Equipay-trips » avec le nombre de lieux,
le temps gagné, le temps de réponse, l'origine des temps (OSRM ou estimation) et le texte
produit, ce qui permet
de suivre la qualité de la rédaction et de préparer un éventuel affinage. Enfin, le trajet
d'une journée peut être ouvert dans l'application Plans d'Apple~\cite{mapkit}, avec
toutes les étapes dans l'ordre et le mode de déplacement choisi.

% ---------------------------------------------------------------------
\subsection{Fonctionnalités transverses}

\subsubsection{Équilibrage des comptes}

À partir des soldes, l'application calcule la liste des remboursements à effectuer
(figure~\ref{fig:code_rembourse}). Les débiteurs et les créanciers sont triés par montant ;
le plus gros débiteur rembourse le plus gros créancier, et l'on recommence avec ce qui
reste. Pour $n$ personnes concernées, cette méthode produit au plus $n - 1$ virements.

\begin{figure}[H]
\begin{lstlisting}[language=Dart]
List<Settlement> settleUp(List<MemberSummary> members) {
  final debtors = [for (final m in members) if (m.balance < -0.005) (m, -m.balance)]
    ..sort((a, b) => b.$2.compareTo(a.$2));
  final creditors = [for (final m in members) if (m.balance > 0.005) (m, m.balance)]
    ..sort((a, b) => b.$2.compareTo(a.$2));
  final result = <Settlement>[];
  var d = 0, c = 0;
  var owe = debtors.isEmpty ? 0.0 : debtors[0].$2;
  var due = creditors.isEmpty ? 0.0 : creditors[0].$2;
  while (d < debtors.length && c < creditors.length) {
    final amount = math.min(owe, due);
    if (amount >= 0.01) result.add(Settlement(debtors[d].$1, creditors[c].$1, (amount * 100).round() / 100));
    owe -= amount;
    due -= amount;
    if (owe < 0.005 && ++d < debtors.length) owe = debtors[d].$2;
    if (due < 0.005 && ++c < creditors.length) due = creditors[c].$2;
  }
  return result;
}
\end{lstlisting}
\caption{Calcul des remboursements qui remettent tous les soldes à zéro}
\label{fig:code_rembourse}
\end{figure}

\subsubsection{Temps réel et mode hors ligne}

Les tables des dépenses, des photos et du programme sont publiées sur le canal temps réel de
Supabase. Lorsqu'un membre ajoute une dépense ou vote, les autres téléphones reçoivent un
événement et rechargent la vue concernée ; les rechargements rapprochés sont regroupés pour
éviter les appels inutiles. Sans réseau, un bandeau l'indique, et l'application se met à jour
d'elle-même au retour de la connexion.

\subsubsection{Rôles, archivage et demandes}

Rejoindre une activité passe par un code de huit caractères, puis par une demande que le
créateur accepte ou refuse. L'archivage est propre à chaque utilisateur : archiver un voyage
terminé ne le masque pas pour les autres. Supprimer une activité est réservé à son créateur ;
un autre membre peut seulement la quitter, à condition de ne pas avoir de dépenses.

% ---------------------------------------------------------------------
\subsection{Extraits de code significatifs et choix d'implémentation}

Plusieurs choix d'implémentation méritent d'être soulignés :

\begin{itemize}
  \item \textbf{La sécurité dans la base, pas dans l'application.} Les droits sont écrits en
        SQL, au plus près des données (figure~\ref{fig:code_rls}). Même une requête construite
        à la main avec un jeton valide ne peut pas lire ou modifier l'activité d'un autre
        groupe.
  \item \textbf{Un repli automatique pour le scanner.} Si le chargement de la version MLX
        échoue, le backend charge la version PyTorch et continue de répondre
        (figure~\ref{fig:code_repli}).
  \item \textbf{Les chiffres ne viennent jamais du modèle quand un calcul est possible.}
        Temps de trajet, ordre de visite, pourcentages et soldes sont calculés ; les modèles
        extraient ou rédigent, ils ne calculent pas.
  \item \textbf{Des services remplaçables.} Chaque service de l'application peut être
        remplacé par une version en mémoire, ce qui a permis d'écrire 62 tests d'interface
        exécutés sans réseau, complétés par 27 tests du backend.
\end{itemize}

\begin{figure}[H]
\begin{lstlisting}[language=SQL2]
create policy trip_items_insert on public.trip_items for insert to authenticated
  with check (
    created_by = (select auth.uid())
    and public.can_edit(activity_id)
    and public.has_program(activity_id)
  );

create policy trip_items_update on public.trip_items for update to authenticated
  using (created_by = (select auth.uid()) or public.is_owner(activity_id))
  with check (public.can_edit(activity_id));
\end{lstlisting}
\caption{Règles d'accès au programme : ajout par un membre, modification par l'auteur ou le créateur}
\label{fig:code_rls}
\end{figure}

\begin{figure}[H]
\begin{lstlisting}[language=Python]
def load(self) -> None:
    """Charge le modèle. En cas d'échec, l'erreur est gardée et l'API répond 503."""
    with self._loading:
        if self.ready:
            return
        engines = ["mlx", "torch"] if config.RECEIPT_ENGINE == "mlx" else ["torch"]
        for engine in engines:
            try:
                (self._load_mlx if engine == "mlx" else self._load_torch)()
                self.engine = engine
                self.error = None
                return
            except Exception as e:
                self.error = f"{type(e).__name__}: {e}"
                log.error("Scanner de reçus (%s) indisponible : %s", engine, self.error)
\end{lstlisting}
\caption{Chargement du scanner avec repli de MLX vers PyTorch (extrait simplifié)}
\label{fig:code_repli}
\end{figure}

% ---------------------------------------------------------------------
\subsection{Exécution locale et choix d'hébergement}

\subsubsection{Exécution sur le Mac}

Pendant le projet, le backend d'IA tourne sur le Mac de développement. Un script de
lancement démarre MLflow, attend qu'il soit prêt, démarre le backend, puis installe et lance
l'application sur l'iPhone ; un second script arrête l'ensemble. Le téléphone joint le
backend par le nom réseau du Mac plutôt que par son adresse IP, qui change d'un réseau à
l'autre. Le téléphone et le Mac doivent simplement être sur le même réseau.

\subsubsection{Comparaison avec un serveur GPU dans le cloud}

\begin{table}[H]
\centering
\small
\begin{tabularx}{\textwidth}{|L{3.0cm}|Y|Y|}
\hline
\rowcolor{entete}\textbf{Critère} & \textbf{Mac local (retenu)} & \textbf{Serveur GPU dans le cloud} \\ \hline
Coût & Nul : matériel déjà disponible & Facturation à l'heure, même sans utilisateur \\ \hline
Mise en place & Immédiate, scripts de lancement & Image, pilotes GPU, HTTPS et nom de domaine à configurer \\ \hline
Disponibilité & Seulement quand le Mac est allumé et sur le même réseau & Accessible partout, à tout moment \\ \hline
Confidentialité & Photos et phrases ne quittent pas le réseau local & Données envoyées à un hébergeur \\ \hline
Mémoire & 18 Go partagés, d'où la quantification en 4 bits & Plus de marge, modèles en 16 bits possibles \\ \hline
Cas d'usage idéal & Développement, démonstration, petit groupe & Publication de l'application au public \\ \hline
\end{tabularx}
\caption{Comparaison entre l'exécution locale et un serveur GPU dans le cloud}
\label{tab:hebergement}
\end{table}

Le tableau~\ref{tab:hebergement} résume ce choix. Pour un projet de démonstration utilisé par
quelques personnes, l'exécution locale offre un coût nul et garde les données sur le réseau
local, au prix d'une disponibilité limitée. Un passage au cloud ne demanderait aucun
changement dans l'application : il suffirait de modifier l'adresse du backend.

% ---------------------------------------------------------------------
\subsection{Conclusion du chapitre}

Ce chapitre a détaillé la réalisation d'\equipay : l'environnement de travail, le choix des
modèles, l'entraînement et l'évaluation de chacun d'eux, l'organisation du trajet, les
fonctionnalités transverses et la stratégie d'exécution retenue. Chaque modèle affiné fait
nettement mieux que son modèle de base sur un jeu de test mesuré de la même façon pour
toutes les versions, et chaque limite connue est compensée dans l'application par une
vérification humaine. Le chapitre suivant présente l'application obtenue à travers ses
principaux écrans.

% =====================================================================
\chapitre{Présentation de l'application réalisée}
% =====================================================================

Après la conception et la réalisation, ce dernier chapitre montre le résultat. Il présente
\equipay telle qu'elle a été testée sur iPhone, à travers un parcours complet : un voyage créé, un membre qui le rejoint, deux dépenses ajoutées avec l'IA, les
comptes équilibrés, puis le programme de la journée organisé. Toutes les captures sont
issues d'une utilisation réelle de l'application, dans la devise de l'activité (le dirham).

% ---------------------------------------------------------------------
\subsection{Aperçu général}

L'écran d'accueil (figure~\ref{fig:ecran_accueil}, à gauche) liste les activités de
l'utilisateur ; le bouton orange permet d'en créer une ou d'en rejoindre une avec son code,
et le mode « Sélectionner » permet d'archiver ou de supprimer plusieurs activités d'un coup.
La barre d'onglets donne accès aux activités, aux demandes pour rejoindre et au profil ; un
quatrième onglet, « Free Card », est réservé à une évolution future. Le profil regroupe les
informations du compte, les activités archivées et les préférences (langue, mode sombre). La
bulle « Support » ouvre l'assistant conversationnel depuis n'importe quel écran principal.

\begin{figure}[H]
\centering
\begin{subfigure}[t]{0.31\textwidth}\centering\capture[0.93]{accueil.jpeg}\caption{Accueil : mes activités}\end{subfigure}\hfill
\begin{subfigure}[t]{0.31\textwidth}\centering\capture[0.93]{profil.jpeg}\caption{Profil de l'utilisateur}\end{subfigure}\hfill
\begin{subfigure}[t]{0.31\textwidth}\centering\capture[0.93]{profil_preferences.jpeg}\caption{Préférences et compte}\end{subfigure}
\caption{Accueil, profil et préférences de l'application}
\label{fig:ecran_accueil}
\end{figure}

% ---------------------------------------------------------------------
\subsection{Page d'une activité}

La page d'une activité (figure~\ref{fig:ecran_activite}) affiche son logo, son type et son
code d'accès, masqué par défaut et copiable en un geste. La carte sarcelle donne le total des
dépenses ; juste en dessous, la carte « Programme » mène au programme du voyage. Les onglets
Dépenses, Membres et Photos donnent accès au détail. Le menu à trois points rassemble les
actions moins fréquentes : modifier, archiver, résumé, graphique des dépenses, statistiques
et suppression. Après l'ajout des deux dépenses du test, le total passe à 65,95 DH et chaque
dépense indique qui l'a payée.

\begin{figure}[H]
\centering
\begin{subfigure}[t]{0.31\textwidth}\centering\capture[0.93]{activite_vide.jpeg}\caption{Activité qui vient d'être créée}\end{subfigure}\hfill
\begin{subfigure}[t]{0.31\textwidth}\centering\capture[0.93]{activite_actions.jpeg}\caption{Menu des actions}\end{subfigure}\hfill
\begin{subfigure}[t]{0.31\textwidth}\centering\capture[0.93]{activite_apres_ajout.jpeg}\caption{Après l'ajout de deux dépenses}\end{subfigure}
\caption{Page d'une activité de type voyage}
\label{fig:ecran_activite}
\end{figure}

% ---------------------------------------------------------------------
\subsection{Rejoindre une activité}

Pour rejoindre un voyage, un ami saisit son code depuis l'accueil. La demande apparaît alors
chez le créateur, dans l'onglet des demandes reçues, avec le nom et le pseudonyme du
demandeur ; le créateur l'accepte ou la refuse (figure~\ref{fig:ecran_demandes}). Le
demandeur suit l'état de ses propres demandes dans l'onglet des demandes envoyées.

\begin{figure}[H]
\centering
\begin{subfigure}[t]{0.31\textwidth}\centering\capture[0.93]{demandes_recues.jpeg}\caption{Demande reçue et acceptée}\end{subfigure}\hspace{0.05\textwidth}
\begin{subfigure}[t]{0.31\textwidth}\centering\capture[0.93]{demandes_envoyees.jpeg}\caption{Demandes envoyées}\end{subfigure}
\caption{Gestion des demandes pour rejoindre une activité}
\label{fig:ecran_demandes}
\end{figure}

% ---------------------------------------------------------------------
\subsection{Ajouter une dépense en scannant un ticket}

Le bouton « Dépense » propose trois façons d'ajouter une dépense : la saisie manuelle, le
scanner de tickets et la description en une phrase (figure~\ref{fig:ecran_scan}). Après la
photo, une animation de balayage indique que l'assistant lit le ticket. Le formulaire
s'ouvre ensuite pré-rempli : sur le ticket du test, l'enseigne (« Shop Name »), le total de
47,95 et la devise ont été lus correctement. Le bandeau « Rempli à partir du ticket »
rappelle de vérifier chaque champ, et l'application signale que la date n'a pas été trouvée
sur le ticket, plutôt que d'en inventer une. Les articles lus sont reportés dans la
description, et la photo du ticket est jointe à la dépense.

\begin{figure}[H]
\centering
\begin{subfigure}[t]{0.235\textwidth}\centering\capture[0.95]{ajout_options.jpeg}\caption{Trois façons d'ajouter}\end{subfigure}\hfill
\begin{subfigure}[t]{0.235\textwidth}\centering\capture[0.95]{scan_lecture.jpeg}\caption{Lecture du ticket}\end{subfigure}\hfill
\begin{subfigure}[t]{0.235\textwidth}\centering\capture[0.95]{scan_formulaire.jpeg}\caption{Formulaire pré-rempli}\end{subfigure}\hfill
\begin{subfigure}[t]{0.235\textwidth}\centering\capture[0.95]{scan_formulaire_suite.jpeg}\caption{Répartition et articles}\end{subfigure}
\caption{Ajout d'une dépense à partir de la photo d'un ticket}
\label{fig:ecran_scan}
\end{figure}

% ---------------------------------------------------------------------
\subsection{Décrire une dépense en une phrase}

La deuxième dépense du test a été saisie en écrivant simplement « Haythame a payé le taxi
18 DH pour nous deux » (figure~\ref{fig:ecran_phrase}). L'écran rappelle les membres de
l'activité et sa devise, puis affiche « L'assistant lit ta phrase… » pendant le traitement.
Le formulaire obtenu contient le titre (Taxi), le montant (18 DH), la devise, la date du jour
(aucune date n'étant donnée dans la phrase), le payeur (Haythame Benladoul) et une
répartition à parts égales entre les deux membres. Sous chaque champ, la mention « Compris
dans ta phrase » indique son origine.

\begin{figure}[H]
\centering
\begin{subfigure}[t]{0.31\textwidth}\centering\capture[0.93]{decrire_saisie.jpeg}\caption{Phrase en cours d'analyse}\end{subfigure}\hfill
\begin{subfigure}[t]{0.31\textwidth}\centering\capture[0.93]{decrire_formulaire.jpeg}\caption{Champs compris}\end{subfigure}\hfill
\begin{subfigure}[t]{0.31\textwidth}\centering\capture[0.93]{decrire_formulaire_suite.jpeg}\caption{Payeur et répartition}\end{subfigure}
\caption{Ajout d'une dépense à partir d'une phrase libre}
\label{fig:ecran_phrase}
\end{figure}

% ---------------------------------------------------------------------
\subsection{Soldes, résumé et graphiques}

Une fois les deux dépenses enregistrées, le résumé (figure~\ref{fig:ecran_soldes}, à gauche)
indique qu'un seul remboursement suffit pour équilibrer les comptes : l'utilisateur a payé
47,95 DH pour une part de 32,98 DH, il doit donc recevoir 14,98 DH, que Haythame lui
rembourse. Le résumé peut être copié pour être envoyé au groupe. Le graphique des dépenses
présente les mêmes données sous trois angles : ce que chacun a payé (73 \% et 27 \%), ce que
chacun a consommé (50 \% chacun) et le solde qui en résulte (+14,98 DH et $-$14,98 DH).

\begin{figure}[H]
\centering
\begin{subfigure}[t]{0.235\textwidth}\centering\capture[0.95]{resume.jpeg}\caption{Résumé}\end{subfigure}\hfill
\begin{subfigure}[t]{0.235\textwidth}\centering\capture[0.95]{graphique_paye.jpeg}\caption{Payé}\end{subfigure}\hfill
\begin{subfigure}[t]{0.235\textwidth}\centering\capture[0.95]{graphique_part.jpeg}\caption{Part}\end{subfigure}\hfill
\begin{subfigure}[t]{0.235\textwidth}\centering\capture[0.95]{graphique_solde.jpeg}\caption{Solde}\end{subfigure}
\caption{Résumé des remboursements et graphiques par participant}
\label{fig:ecran_soldes}
\end{figure}

La page de statistiques (figure~\ref{fig:ecran_stats}) complète cette vue : total, nombre de
dépenses, moyenne par dépense et par participant, jour le plus dépensier, situation
personnelle de l'utilisateur, faits marquants (plus grosse dépense, plus gros payeur) et
répartition des dépenses par jour de la semaine.

\begin{figure}[H]
\centering
\begin{subfigure}[t]{0.31\textwidth}\centering\capture[0.93]{insights_1.jpeg}\caption{Indicateurs et situation}\end{subfigure}\hspace{0.05\textwidth}
\begin{subfigure}[t]{0.31\textwidth}\centering\capture[0.93]{insights_2.jpeg}\caption{Faits marquants}\end{subfigure}
\caption{Statistiques des dépenses d'une activité}
\label{fig:ecran_stats}
\end{figure}

% ---------------------------------------------------------------------
\subsection{Programme du voyage}

Pour le test, trois visites ont été ajoutées au programme : les Champs-Élysées, Notre-Dame
et la place Charles-de-Gaulle, le mercredi 14 octobre 2026 à 10 h, 12 h et 15 h
(figure~\ref{fig:ecran_programme}). Le formulaire demande un nom, un lieu, une adresse, une
date et une heure ; le lieu et l'adresse servent à placer automatiquement la visite sur la
carte. Chaque proposition peut recevoir des votes pour ou contre, affichés en direct. Le
programme se consulte en liste, en calendrier (une frise par jour, dans l'ordre des heures)
ou sur la carte Apple Plans, où chaque visite est un repère numéroté dans l'ordre
chronologique et coloré selon les votes ; un filtre permet de n'afficher qu'un jour.

\begin{figure}[H]
\centering
\begin{subfigure}[t]{0.235\textwidth}\centering\capture[0.95]{programme_ajout.jpeg}\caption{Ajout d'une visite}\end{subfigure}\hfill
\begin{subfigure}[t]{0.235\textwidth}\centering\capture[0.95]{programme_liste.jpeg}\caption{Liste et votes}\end{subfigure}\hfill
\begin{subfigure}[t]{0.235\textwidth}\centering\capture[0.95]{programme_calendrier.jpeg}\caption{Calendrier}\end{subfigure}\hfill
\begin{subfigure}[t]{0.235\textwidth}\centering\capture[0.95]{programme_carte.jpeg}\caption{Carte Apple Plans}\end{subfigure}
\caption{Programme d'un voyage : ajout, liste, calendrier et carte}
\label{fig:ecran_programme}
\end{figure}

% ---------------------------------------------------------------------
\subsection{Organisation du trajet}

Le bouton « Organiser le trajet » ouvre les options (figure~\ref{fig:ecran_trajet}) : le
mode de déplacement (à pied, à vélo, en voiture), le nombre maximal de visites par jour et
une consigne facultative. Sur le test, le trajet à pied a été demandé. Les heures des trois
visites étant déjà fixées et l'ordre déjà optimal, l'application l'indique (« Votre ordre est
déjà le meilleur ») et affiche le temps de marche total, soit 2 h 00. Pour chaque étape, elle
donne l'heure, puis le temps et la distance jusqu'à l'étape suivante (52 min et 3,9 km, puis
1 h 08 et 5,1 km), et trace l'itinéraire dans les rues sur la carte. Le bouton « Ouvrir
l'itinéraire dans Plans » lance la navigation dans l'application Plans avec toutes les
étapes, et le créateur peut enregistrer le programme proposé avec « Appliquer ce
programme ».

\begin{figure}[H]
\centering
\begin{subfigure}[t]{0.31\textwidth}\centering\capture[0.93]{trajet_options.jpeg}\caption{Options du trajet}\end{subfigure}\hspace{0.05\textwidth}
\begin{subfigure}[t]{0.31\textwidth}\centering\capture[0.93]{trajet_resultat.jpeg}\caption{Trajet proposé}\end{subfigure}
\caption{Organisation automatique du trajet d'une journée}
\label{fig:ecran_trajet}
\end{figure}

% ---------------------------------------------------------------------
\subsection{Conclusion du chapitre}

Les tests menés sur iPhone couvrent l'ensemble du parcours de l'application : création d'une
activité et arrivée d'un membre, ajout de dépenses par le scanner et par une phrase,
équilibrage des comptes, statistiques, programme partagé et organisation du trajet. Chaque
fonction d'IA a été vérifiée en conditions réelles, et chaque fois le résultat a été proposé à
l'utilisateur avant d'être enregistré. Ces tests restent menés à petite échelle (un voyage,
deux membres, quelques dépenses) : le comportement de l'application avec de grands groupes
et un usage prolongé reste à observer.

% ---------------------------------------------------------------------
%  Conclusion générale
% ---------------------------------------------------------------------
\pagetitre{CONCLUSION GÉNÉRALE}

Ce projet nous a permis de concevoir et de réaliser \equipay, une application mobile de
partage des dépenses de groupe qui fait de l'intelligence artificielle un outil concret
plutôt qu'une vitrine. Partant d'un constat simple (la saisie manuelle décourage les
utilisateurs, et rien ne relie les dépenses d'un voyage à son organisation), nous avons
construit une application complète : Flutter pour le mobile, Supabase pour les données
partagées et sécurisées, un backend FastAPI pour les modèles, et MLflow pour suivre leur
cycle de vie.

Sur le plan de l'IA, quatre fonctions ont été intégrées. L'assistant conversationnel, affiné
à partir de Qwen3-4B sur un jeu de données que nous avons construit, passe d'un score de
0,276 à 0,770 et a appris à refuser les questions hors sujet. Le scanner de tickets, affiné à
partir de Qwen3-VL-4B, trouve le bon total sur 97,2 \% des tickets de test et sur la totalité
des tickets en dirhams et en euros. La saisie par phrase atteint un score de 0,976. Enfin,
l'organisation du trajet combine un calcul d'itinéraire, un algorithme d'ordre de visite et
un modèle de langage chargé uniquement de la rédaction. La conversion en 4 bits avec MLX a
réduit de moitié la mémoire nécessaire, ce qui permet de servir tous ces modèles sur un
simple ordinateur portable.

Au-delà des résultats chiffrés, ce projet nous a appris plusieurs leçons. La première
version du scanner, excellente sur son jeu de test mais inutilisable sur nos propres
tickets, nous a montré qu'un bon score ne vaut que si le jeu de test ressemble à l'usage
réel. La séparation entre ce que calcule un algorithme et ce que rédige un modèle s'est
révélée essentielle pour obtenir des résultats fiables. Et la démarche MLOps (données
versionnées, évaluations comparables, choix d'un champion) a transformé une suite d'essais en
une progression mesurable. Le travail en binôme, avec relecture croisée et décisions prises
à deux, a été un apprentissage directement transposable en entreprise.

L'application reste un prototype validé à petite échelle. Les perspectives sont claires :
héberger le backend sur un serveur GPU pour rendre l'IA disponible partout, publier
l'application sur les magasins d'applications (ce qui imposera la connexion avec un compte
Apple), compléter la traduction en anglais et en arabe déjà prévue dans les préférences,
enrichir le jeu de données de l'assistant pour réduire ses réponses inventées sur les
questions de support, et ajouter les notifications. Au terme de ce projet, l'objectif initial
est atteint : \equipay montre qu'il est possible d'intégrer des modèles de langage affinés
dans une application du quotidien, de façon mesurée, traçable et utile.


% =====================================================================
%  Références : seules les sources citées dans le texte apparaissent
% =====================================================================
\bibliographystyle{unsrt-fr}
\bibliography{references}

\end{document}
